\documentclass[11pt]{article}
\usepackage[T1]{fontenc}
\usepackage{amsmath,amssymb,amsthm,mathtools}
\usepackage[margin=1in]{geometry}
\usepackage[hidelinks]{hyperref}

\newtheorem{proposition}{Proposition}
\newtheorem{theorem}[proposition]{Theorem}
\newtheorem{lemma}[proposition]{Lemma}
\newtheorem{corollary}[proposition]{Corollary}
\newtheorem{conjecture}[proposition]{Conjecture}
\newtheorem{definition}[proposition]{Definition}
\newtheorem{remark}[proposition]{Remark}
\newtheorem{example}[proposition]{Example}
\newcommand{\G}{\mathfrak{G}}
\newcommand{\Q}{\mathcal{Q}}
\newcommand{\Fix}{\operatorname{Fix}}
\newcommand{\Left}{\operatorname{Left}}
\newcommand{\Out}{\operatorname{Out}}
\newcommand{\ellength}{\ell}
\newcommand{\qdeg}{\operatorname{qdeg}}
\DeclareMathOperator{\supp}{supp}

\title{Triple quantum K-theory of flag varieties}
\author{Matt Samuel}
\date{\today}

\begin{document}
\maketitle

\begin{abstract}
We prove a triple quantum K-Pieri formula expanding the product of a
quantum double Grothendieck polynomial and a quantized top block. Its
coefficients are indexed by endpoints of admissible chains in the
quantum Bruhat graph. A cyclic rotation of values turns these chains
into the Pieri chains of Lenart and Sottile in the classical
$k$-Bruhat order; this gives uniqueness of the chain to each endpoint
and a closed formula for its quantum degree. We prove the formula for
$k=1$, at $y=0$ for all $k$, and on the hyperplane $z=1/\beta$ for all
$k$, where it is an equivariant quantum $K$-Chevalley formula for the
line bundle $\mathcal O(-\varpi_k)$ indexed by chains; the general
equivariant case is reduced to a junction identity whose remaining
part is stated precisely.
\end{abstract}

\section{Polynomial conventions}

Let $S_\infty$ denote the permutations of the positive integers with
finite support. We use one-line notation, let $\ellength(w)$ be the
inversion number of $w$, and write $t_{ab}$ for the transposition
interchanging positions $a<b$. Thus $wt_{ab}$ is obtained from $w$
by interchanging $w(a)$ and $w(b)$.

Let $\beta$, $z$, and the entries of $x=(x_1,x_2,\ldots)$,
$y=(y_1,y_2,\ldots)$, and $q=(q_1,q_2,\ldots)$ be independent
indeterminates. Put
\[
 X_i=1+\beta x_i,\qquad Q_i=\beta^2q_i,\qquad
 \overline y_a=-\frac{y_a}{1+\beta y_a}.
\]
Our formal group law is $a\oplus b=a+b+\beta ab$, so that
$a\oplus\overline a=0$. In particular,
$z+\overline y_a$ below denotes an \emph{ordinary} sum, not
$z\oplus\overline y_a$.

We normalize the double Grothendieck polynomials by
\[
 \G_{w_0^{(n)}}(x;y)
   =\prod_{\substack{i,j\geq 1\\i+j\leq n}}(x_i\oplus y_j),
\]
and the isobaric divided-difference recursion
\[
 \G_{ws_i}(x;y)=\pi_i^\beta\G_w(x;y)
 \quad\text{if }\ellength(ws_i)<\ellength(w),\qquad
 \pi_i^\beta f
 =\frac{(1+\beta x_{i+1})f-(1+\beta x_i)s_i f}
        {x_i-x_{i+1}}.
\]
Here $s_i$ interchanges $x_i,x_{i+1}$ and fixes all other variables.
In this convention $\G_{21}=x_1+y_1+\beta x_1y_1$, and
$\G_w|_{\beta=0}$ is the double Schubert polynomial with second
alphabet $-y$.

For $0\leq r\leq j$ define
\begin{equation}\label{eq:quantum-elementary}
 F_r^j(X;q,\beta)
 =\sum_{\substack{J\subseteq\{1,\ldots,j\}\\|J|=r}}
    \left(\prod_{a\in J}X_a\right)
    \left(\prod_{\substack{a\in J\\a+1\notin J}}
                     (1-\beta^2q_a)\right).
\end{equation}
In particular $F_0^j=1$.
The Lenart--Maeno quantization map $\Q$ is linear over the coefficient
parameters and is specified on standard elementary monomials by
\begin{equation}\label{eq:quantization}
 \Q\left(\prod_{j=1}^{N-1}e_{r_j}(X_1,\ldots,X_j)\right)
   =\prod_{j=1}^{N-1}F_{r_j}^j(X;q,\beta),
 \qquad 0\leq r_j\leq j.
\end{equation}
One takes $N$ sufficiently large for the polynomial being quantized.
The standard elementary monomials form a basis of the space with
$\deg_{X_i}\leq N-i$, $1\leq i<N$; this specifies the map on that
space. The map is not asserted to be an algebra homomorphism.
We set
\[
 \G_u^q(x;y)=\Q\bigl(\G_u(x;y)\bigr).
\]
At $\beta=-1$ this is the usual quantum double Grothendieck normalization
associated with Lenart--Maeno quantization
\cite{naito2024pieri,maeno2024presentation2}.

The top block appearing in our formula is
\begin{equation}\label{eq:top-block}
 B_k^q(x;z)
 :=\Q\left(\prod_{i=1}^{k}(x_i+z)\right)
 =\beta^{-k}\sum_{r=0}^{k}
       \bigl(-(1-\beta z)\bigr)^{k-r}F_r^k(X;q,\beta).
\end{equation}
The apparent negative powers of $\beta$ in this expression cancel.
At $\beta=0$ the block is understood by polynomial specialization,
not by separately substituting into $\beta^{-k}$.
We set $B_0^q=1$.
Notice that \eqref{eq:top-block} quantizes the ordinary factors
$x_i+z$, not the formal-group factors $x_i\oplus z$.

\section{Chains and endpoint factors}

Fix $u\in S_\infty$ and $k\geq 1$. Choose a finite rank $R>k$ with
$u\in S_R$. The directed quantum Bruhat graph on $S_R$ has an edge
\[
 v\xrightarrow{(a,b)}vt_{ab}\qquad(1\leq a<b\leq R)
\]
when either
\begin{align*}
 \ellength(vt_{ab})&=\ellength(v)+1
       &&\text{(a Bruhat edge)},\\
 \ellength(vt_{ab})&=\ellength(v)-2(b-a)+1
       &&\text{(a quantum edge)}.
\end{align*}
A Bruhat edge has weight $1$; a quantum edge labeled $(a,b)$ has
weight $q_aq_{a+1}\cdots q_{b-1}$.
Order labels by
\[
 (a,b)\prec(c,d)
 \quad\Longleftrightarrow\quad
 b>d\ \text{ or }\ (b=d\text{ and }a<c).
\]

\begin{definition}\label{def:chain}
An admissible $k$-chain from $u$ in $S_R$ is a directed path
\[
 \gamma:\quad
 u=v_0\xrightarrow{(a_1,b_1)}v_1
       \xrightarrow{(a_2,b_2)}\cdots
       \xrightarrow{(a_s,b_s)}v_s=w
\]
in the quantum Bruhat graph satisfying:
\begin{enumerate}
 \item $a_i\leq k<b_i$ for every $i$;
 \item all labels $(a_i,b_i)$ are distinct;
 \item $b_1\geq b_2\geq\cdots\geq b_s$;
 \item if $i<s$ and $a_i=a_j$ for some $j<i$, then
       $(a_i,b_i)\prec(a_{i+1},b_{i+1})$.
\end{enumerate}
The empty path, with endpoint $u$, is included.
Let $E_R(u,k)$ be the set of endpoints of these paths.
Write $s(\gamma)$ for the number of edges, and define
$D(\gamma)=(D_1,\ldots,D_{R-1})$ by
\[
 D_j(\gamma)
 =\#\{\,i:\text{the $i$th edge is quantum and }a_i\leq j<b_i\,\}.
\]
Thus the weight of $\gamma$ is
$q^{D(\gamma)}=\prod_{j=1}^{R-1}q_j^{D_j(\gamma)}$.
\end{definition}

 Directly from the edge definitions, every such path satisfies
 \begin{equation}\label{eq:length-degree}
  s(\gamma)=\ellength(w)-\ellength(u)
                +2\sum_{j=1}^{R-1}D_j(\gamma).
 \end{equation}

 For an endpoint $w$ define three disjoint sets of \emph{values}
 from the initial window $u(1),\ldots,u(k)$:
\begin{align*}
 \Fix(u,w;k)
   &=\{u(i):1\leq i\leq k,\ w(i)=u(i)\},\\
 \Left(u,w;k)
   &=\{u(i):1\leq i\leq k,\
                      w^{-1}(u(i))<i\},\\
 \Out(u,w;k)
   &=\{u(i):1\leq i\leq k\}
                        \setminus\bigl(\Fix(u,w;k)\cup\Left(u,w;k)\bigr).
\end{align*}
Set
\[
 m(u,w;k)=|\Left(u,w;k)|+|\Out(u,w;k)|
          =k-|\Fix(u,w;k)|.
\]
The name $\Out$ includes values that move right while remaining in
the window, as well as values that leave the window.

 \section{The triple quantum K-Pieri formula}

 We state the triple quantum K-Pieri formula, expanding a quantum
 double Grothendieck polynomial in the $y$-equivariant basis times a
 quantized top block carrying a second parameter $z$.

 \begin{theorem}[Triple quantum K-Pieri formula]
\label{prop:top}
 Let $k\geq1$ and $u\in S_\infty$. Set
 \[
  R_0=\max\{\max(\{2\}\cup\supp(u)),\,k+1\}+k,
  \qquad E=E_{R_0}(u,k).
 \]
 For each endpoint $w\in E$, all admissible chains from $u$ to $w$
 have the same quantum degree $D(u,w;k)$ and length $s(u,w;k)$.
 Moreover, $s(u,w;k)\geq m(u,w;k)$, and
 \begin{equation}\label{eq:main}
   B_k^q(x;z)\,\G_u^q(x;y)
     =\sum_{w\in E}
        C_{u,w}^{(k)}(y,z,q,\beta)\,\G_w^q(x;y),
 \end{equation}
 where
 \begin{equation}\label{eq:coefficient}
 \boxed{\begin{aligned}
 C_{u,w}^{(k)}
   ={}&q^{D(u,w;k)}\beta^{s(u,w;k)-m(u,w;k)}
       \prod_{a\in\Fix(u,w;k)}
          \left(z-\frac{y_a}{1+\beta y_a}\right)\\
     &\quad{}\times(1-\beta z)^{|\Left(u,w;k)|}
       \prod_{a\in\Out(u,w;k)}(1+\beta y_a)^{-1}.
 \end{aligned}}
 \end{equation}
 Here $B_k^q(x;z)=\Q(\prod_{i=1}^k(x_i+z))$ is defined in
 \eqref{eq:top-block}; $D$ records the quantum edges in a chain;
 $s$ is its number of edges; and
 $m=|\Left|+|\Out|=k-|\Fix|$. The sets $\Fix,\Left,\Out$ record
  whether each value among $u(1),\ldots,u(k)$ stays fixed, moves to an
  earlier position, or otherwise moves (to a later position or out of
  the first $k$ positions). The admissible chains and their quantum
  weights are defined above.
 \end{theorem}

 The formula is indexed by endpoints, so each endpoint appears once.
 The chain-independence of $D$ and $s$ is part of the theorem; it is
 proved in Corollary~\ref{cor:rotation}, which shows that there is
 exactly one admissible chain to each endpoint and gives
 $D$ and $s$ in closed form, see \eqref{eq:D-closed}.

 \begin{remark}[Rank and polynomial setting]
 The rank $R_0$ just makes the endpoint set finite; it does not impose
 relations from a finite flag variety. For $k=0$ the formula reduces
 to $\G_u^q=\G_u^q$.
 \end{remark}

 \begin{remark}[A useful degree check]
 Let $|D|=\sum_jD_j$. By \eqref{eq:length-degree}, the prefactor in
 \eqref{eq:coefficient} can also be written as
 \[
 q^D\beta^{s-m}
   =Q^D\beta^{\ellength(w)-\ellength(u)-m}.
\]
The first expression is preferable for specialization at $\beta=0$,
since the inequality $s\geq m$ makes nonnegativity of the
power explicit. With
$\deg x_i=\deg y_i=\deg z=1$, $\deg\beta=-1$, and $\deg q_i=2$,
the coefficient has degree $k+\ellength(u)-\ellength(w)$.
\end{remark}

\section{Proof calculations}
\label{sec:proof}

The arguments below establish the rank bound, the uniqueness of the
admissible chain to each endpoint together with closed formulas for
$s$ and $D$ (by a rotation of values to the classical $k$-Bruhat
order), polynomiality of the
top block, the full formula for $k=1$, the full formula at $y=0$
for all $k$ (Theorem~\ref{thm:noneq}, from the Naito--Sagaki
marked-chain formula), and the full equivariant formula at
$z=1/\beta$ for all $k$ (Theorem~\ref{thm:a-zero}, by the junction
identity). The remaining step for the general equivariant
formula is specified at the end of the section; that step is not
proved here.

\subsection{Finiteness and the power of \texorpdfstring{$\beta$}{beta}}

\begin{proposition}\label{prop:rank}
Put $n(u)=\max(\{1\}\cup\supp(u))$ and
$M=\max\{n(u),k\}+1$.
Every label in an admissible $k$-chain from $u$ has second
entry at most $M$. Consequently its endpoint set stabilizes
in every rank containing $u$ and $\{1,\ldots,M\}$.
For every such chain $\gamma$ with endpoint $w$,
\[
 s(\gamma)\geq m(u,w;k).
\]
\end{proposition}

\begin{proof}
The first label has maximal second entry, by the ordering condition.
Suppose it were $(a,b)$ with $b>M$. Then $a\leq k<b-1$,
$u(b)=b$, $u(b-1)=b-1$, and $u(a)<b-1$.
The intermediate value $u(b-1)$ prevents $(a,b)$ from giving a
Bruhat cover. Since $u(a)<u(b)$, the edge cannot be quantum either.
This contradiction proves the rank bound.
Adding fixed positions after $M$ changes neither the relevant
inversion differences nor the permitted paths, proving stabilization.

Every position $i\leq k$ with $w(i)\ne u(i)$ must occur in at least
one transposition of the path. Each transposition involves exactly
one position in $\{1,\ldots,k\}$. There are
$k-|\Fix(u,w;k)|=m(u,w;k)$ changed positions, so at least this
many edges are required.
\end{proof}

This proves that the displayed power $\beta^{s-m}$ is nonnegative
for each path, without presupposing independence of the path.
Independence of $s$ and $D$ for a fixed endpoint, and in fact
uniqueness of the path, is proved next.

\subsection{Rotation to the classical \texorpdfstring{$k$}{k}-Bruhat order}
\label{sec:rotation}

The quantum Bruhat graph has a cyclic symmetry, which we derive from
the following reformulation of the edge condition. For $v\in S_R$ and
$1\leq a<b\leq R$ put
\[
 A(v;a,b)=\{v(r):a<r<b\},
\]
the set of values at the positions strictly between $a$ and $b$.
For distinct $x,y\in\{1,\ldots,R\}$ let $\operatorname{arc}(x,y)$ be
the set of values met strictly between $x$ and $y$ when moving
upward cyclically from $x$ to $y$:
\[
 \operatorname{arc}(x,y)=
 \begin{cases}
  \{x+1,\ldots,y-1\}& x<y,\\
  \{x+1,\ldots,R\}\cup\{1,\ldots,y-1\}& x>y.
 \end{cases}
\]

\begin{lemma}[Cyclic form of the edge condition]\label{lem:arc}
Let $v\in S_R$ and $a<b$. Then $v\xrightarrow{(a,b)}vt_{ab}$ is an
edge of the quantum Bruhat graph if and only if
\[
 A(v;a,b)\cap\operatorname{arc}(v(a),v(b))=\varnothing .
\]
The edge is a Bruhat edge if $v(a)<v(b)$ and a quantum edge if
$v(a)>v(b)$.
\end{lemma}

\begin{proof}
Passing from $v$ to $vt_{ab}$ changes the inversion status of the pair
of positions $(a,b)$ and of the pairs $(a,r)$, $(r,b)$ with $a<r<b$
exactly when $v(r)$ lies strictly between $v(a)$ and $v(b)$; no other
pair is affected. Hence, if $v(a)<v(b)$,
\[
 \ellength(vt_{ab})=\ellength(v)+1+2\,\#\{r:a<r<b,\ v(a)<v(r)<v(b)\},
\]
and if $v(a)>v(b)$,
\[
 \ellength(vt_{ab})=\ellength(v)-1-2\,\#\{r:a<r<b,\ v(b)<v(r)<v(a)\}.
\]
A Bruhat edge therefore means $v(a)<v(b)$ and
$A(v;a,b)\cap(v(a),v(b))=\varnothing$. A quantum edge means
$v(a)>v(b)$ and that all $b-a-1$ intermediate values lie in
$(v(b),v(a))$; since $A(v;a,b)$ does not contain $v(a)$ or $v(b)$,
this says that $A(v;a,b)$ avoids the complement of $[v(b),v(a)]$,
which is $\operatorname{arc}(v(a),v(b))$.
\end{proof}

For $0\leq\vartheta\leq R$ let $\rho=\rho_\vartheta$ be the cyclic
shift of $\{1,\ldots,R\}$ given by
\[
 \rho(x)=x-\vartheta\ \ (x>\vartheta),\qquad
 \rho(x)=x-\vartheta+R\ \ (x\leq\vartheta),
\]
so that $\rho(\vartheta)=R$ is the largest value when
$\vartheta\geq1$, and the relabelled order of values is
$\vartheta+1<\vartheta+2<\cdots<R<1<\cdots<\vartheta$.
For $v\in S_R$ write $\rho v=\rho\circ v$ for the permutation
obtained by relabelling the values of $v$. Since $\rho$ is a cyclic
shift, $\rho(\operatorname{arc}(x,y))=\operatorname{arc}(\rho x,\rho y)$,
and clearly $A(\rho v;a,b)=\rho(A(v;a,b))$. Lemma~\ref{lem:arc} gives
at once:

\begin{corollary}[Cyclic symmetry]\label{cor:cyclic}
$(a,b)$ is an edge from $v$ if and only if it is an edge from
$\rho v$. The two edges have the same type unless exactly one of
$v(a),v(b)$ is at most $\vartheta$: the edge from $v$ is quantum and
the edge from $\rho v$ is Bruhat exactly when
$v(b)\leq\vartheta<v(a)$, and the edge from $v$ is Bruhat and the edge
from $\rho v$ is quantum exactly when $v(a)\leq\vartheta<v(b)$.
\end{corollary}

The point of the following lemma is that, for the paths we consider,
a single rotation removes all quantum edges at once. It uses only the
distinctness of the labels and the weak decrease of the upper
indices, not condition~4 of Definition~\ref{def:chain}.

\begin{lemma}[Rotation lemma]\label{lem:rotation}
Let $u\in S_R$, $1\leq k<R$, $\vartheta=u(k+1)$ and $\rho=\rho_\vartheta$.
Let
\[
 u=v_0\xrightarrow{(a_1,b_1)}v_1\xrightarrow{(a_2,b_2)}\cdots
   \xrightarrow{(a_s,b_s)}v_s
\]
be a directed path in the quantum Bruhat graph whose labels satisfy
$a_i\leq k<b_i$, are pairwise distinct, and have
$b_1\geq b_2\geq\cdots\geq b_s$. Then every edge of the rotated path
$\rho v_0\to\rho v_1\to\cdots\to\rho v_s$ is a Bruhat edge.
Consequently the $i$th edge of the original path is quantum if and
only if $v_{i-1}(b_i)\leq\vartheta<v_{i-1}(a_i)$.
\end{lemma}

\begin{proof}
Put $v_i'=\rho v_i$ and $N=R=\rho(\vartheta)$, the largest value;
$N$ sits at position $k+1$ in $v_0'$. By Corollary~\ref{cor:cyclic}
the rotated path is a directed path with the same labels, and the
last assertion follows from the first by the same corollary. We show
that no rotated edge is quantum.

Position $k+1$ is involved in a label $(a,b)$ only when $b=k+1$.
Let $(c_1,k+1),\ldots,(c_m,k+1)$ be the labels with upper index
$k+1$, in order; by the weak decrease of upper indices they are the
last $m$ labels of the path, and by distinctness the $c_l$ are
pairwise distinct. Thus $N$ occupies position $k+1$ until the label
$(c_1,k+1)$ is applied, and position $c_1$ from then on.

Consider first an edge with label $(a,b)$, $b>k+1$. Then
$a<k+1<b$, so $N\in A(v';a,b)$ for the permutation $v'$ at which it
is applied. A quantum edge from $v'$ would have
$A(v';a,b)\subseteq(v'(b),v'(a))$, which cannot contain the largest
value. Hence the edge is Bruhat. The edge $(c_1,k+1)$ is Bruhat
because the value $N$ at position $k+1$ exceeds the value at $c_1$.

It remains to treat $(c_l,k+1)$ for $l\geq2$. We argue by induction
on $l$, assuming that $(c_1,k+1),\ldots,(c_{l-1},k+1)$ are Bruhat
edges of the rotated path. Put $e_0=N$ and let $e_h$ be the value at
position $c_h$ immediately before the edge $(c_h,k+1)$, for
$1\leq h\leq l$. Immediately before $(c_h,k+1)$ the value at position
$k+1$ is $e_{h-1}$, so the Bruhat property of these edges gives
\[
 N=e_0>e_1>\cdots>e_{l-1},
\]
and after $(c_h,k+1)$ the value $e_{h-1}$ sits at position $c_h$,
where it stays, since $c_h$ is not touched again. Position $c_l$ is
touched by none of the labels $(c_h,k+1)$, $h<l$, so $x:=e_l$ has been
the value at position $c_l$ throughout. For $h<l$ with $c_h<c_l$ the
position $c_l$ lies strictly between $c_h$ and $k+1$, and the Bruhat
condition for the edge $(c_h,k+1)$ gives $x\notin(e_h,e_{h-1})$.

Suppose that the edge $(c_l,k+1)$ were quantum. Then $x>e_{l-1}$ and
every value at a position strictly between $c_l$ and $k+1$ lies in
$(e_{l-1},x)$; in particular $e_{h-1}<x$ whenever $h<l$ and $c_h>c_l$.
Since $e_{l-1}<x<e_0$ and $x$ differs from every $e_h$, there is a
unique $h\in\{1,\ldots,l-1\}$ with $e_h<x<e_{h-1}$. If $c_h<c_l$ this
contradicts $x\notin(e_h,e_{h-1})$; if $c_h>c_l$ it contradicts
$e_{h-1}<x$. Hence the edge is Bruhat.
\end{proof}

A saturated chain in the $k$-Bruhat order is a path of Bruhat edges
with labels $(a,b)$, $a\leq k<b$. Lenart and Sottile
\cite{lenart2007pieri} consider such chains subject to two conditions,
(P1): $b_1\geq\cdots\geq b_s$, and (P0): if $i<s$ and $a_i=a_j$ for
some $j<i$, then $(a_i,b_i)\prec(a_{i+1},b_{i+1})$, with the same
order $\prec$ on labels as above; these are conditions 3 and 4 of
Definition~\ref{def:chain}. Distinctness of labels is automatic for
Bruhat chains satisfying (P1), since a repeated label would force
equal upper indices in between and a return to a smaller permutation.
Their uniqueness theorem \cite[Theorem~2.2]{lenart2007pieri} states
that for given $v,w$ there is at most one saturated chain in the
$k$-Bruhat order from $v$ to $w$ satisfying (P0) and (P1).

\begin{corollary}[Classical reduction of the chains]\label{cor:rotation}
Let $u\in S_R$, $1\leq k<R$ with $R\geq\max\{n(u),k\}+1$,
$\vartheta=u(k+1)$ and $\rho=\rho_\vartheta$. Applying $\rho$ to all
vertices is a bijection from the admissible $k$-chains from $u$ onto
the saturated chains in the $k$-Bruhat order starting at $\rho u$ that
satisfy (P0) and (P1); labels are preserved. Consequently:
\begin{enumerate}
\item For every $w\in E(u,k)$ there is exactly one admissible
 $k$-chain from $u$ to $w$. In particular $s(u,w;k)$ and
 $D(u,w;k)$ are well defined, and the chain is a shortest directed
 path from $u$ to $w$ in the quantum Bruhat graph.
\item Along an admissible chain the value at a position $i\leq k$
 changes exactly at the labels with lower index $i$, and in the
 rotated order it strictly increases at each of them; the value at a
 position $b>k$ changes exactly at the labels with upper index $b$
 and strictly decreases in the rotated order. In particular
 $w(i)\neq u(i)$ as soon as $i$ is the lower index of some label, so
 $m(u,w;k)$ equals the number of distinct lower indices of the chain.
\item Each position is involved in at most one quantum edge, and
\begin{equation}\label{eq:D-closed}
 D_j(u,w;k)=
 \begin{cases}
  \#\{\,i\leq j:\ u(i)>u(k+1)\geq w(i)\,\}& j<k,\\[2pt]
  \#\{\,b>j:\ u(b)\leq u(k+1)<w(b)\,\}& j\geq k,
 \end{cases}
 \qquad
 s(u,w;k)=\ellength(w)-\ellength(u)+2\sum_jD_j(u,w;k).
\end{equation}
\end{enumerate}
\end{corollary}

\begin{proof}
By Lemma~\ref{lem:rotation} the rotation of an admissible $k$-chain
is a path of Bruhat edges with the same labels, hence a saturated
chain in the $k$-Bruhat order, and conditions 3 and 4 are (P1) and
(P0). Conversely, by Corollary~\ref{cor:cyclic} a chain in the
$k$-Bruhat order from $\rho u$ with (P0) and (P1) rotates back to a
directed path from $u$ with the same labels, which are distinct by the
remark above; so it is an admissible $k$-chain. The inverse rotation
is $\rho^{-1}$, so this is a bijection.

(1) Two admissible chains from $u$ to $w$ rotate to two chains from
$\rho u$ to $\rho w$ satisfying (P0) and (P1), which coincide by the
uniqueness theorem of Lenart and Sottile. Each edge of the quantum
Bruhat graph raises the length by at most one, so a directed path from
$u$ to $w$ has at least $\ellength(\rho w)-\ellength(\rho u)$ edges;
the rotated chain has exactly this many.

(2) A Bruhat edge $(a,b)$ replaces the value at $a$ by the larger
value at $b$ and the value at $b$ by the smaller value at $a$; no
other position changes.

(3) By Lemma~\ref{lem:rotation} the $i$th edge is quantum exactly
when $v_{i-1}(b_i)\leq\vartheta<v_{i-1}(a_i)$, that is, when the value
entering position $a_i$ is at most $\vartheta$ while the value leaving
it exceeds $\vartheta$, and the value entering position $b_i$ exceeds
$\vartheta$ while the value leaving it is at most $\vartheta$. In the
rotated order the values $>\vartheta$ precede the values
$\leq\vartheta$, so by (2) this transition happens at most once at
each position, namely when the initial value at $a_i$ exceeds
$\vartheta$ and the final one does not, and when the initial value at
$b_i$ is at most $\vartheta$ and the final one exceeds it. A quantum
edge $(a,b)$ contributes $1$ to $D_j$ exactly when $a\leq j<b$. For
$j<k$ every quantum edge has $b>k>j$, so $D_j$ counts the quantum
edges with $a\leq j$, which by the above is the first formula. For
$j\geq k$ every quantum edge has $a\leq k\leq j$, so $D_j$ counts the
quantum edges with $b>j$, giving the second formula. The formula for
$s$ is \eqref{eq:length-degree}.
\end{proof}

\begin{remark}
The corollary identifies the indexing set of Theorem~\ref{prop:top}:
$w\in E(u,k)$ if and only if $\rho w$ is reached from $\rho u$ by a
Lenart--Sottile chain, that is, $\rho u\xrightarrow{c(k)}\rho w$ in
the notation of \cite{lenart2007pieri}, where $\rho$ rotates the
values so that $u(k+1)$ becomes the largest. Thus the quantum
endpoint set for $u$ is the classical endpoint set for $\rho u$, and
the quantum degree records which window values drop from above
$u(k+1)$ to at most $u(k+1)$, and which outside values rise across
it. We have checked \eqref{eq:D-closed} against the direct
enumeration of chains for all $u\in S_n$, $n\leq7$, and all $k$.
Lemma~\ref{lem:rotation} also shows that the chains, but not the line
paths of Section~\ref{sec:line-chevalley} below, can be analysed in
the classical $k$-Bruhat order; backtracks such as
$(k-1,k),(k-1,k)$ are inherently non-classical.
\end{remark}

\subsection{An explicit polynomial top block}

Write $a=1-\beta z$. For $S\subseteq\{1,\ldots,k\}$ with no two
consecutive elements, put
\[
 \epsilon(S)=
 \begin{cases}1&k\in S,\\0&k\notin S,\end{cases}
 \qquad
 U(S)=\{1,\ldots,k\}\setminus\bigl(S\cup(S+1)\bigr),
 \quad S+1=\{j+1:j\in S\}.
\]

\begin{proposition}\label{prop:block-polynomial}
The top block has the polynomial expression
\begin{equation}\label{eq:block-matchings}
 B_k^q(x;z)=
 \sum_{\substack{S\subseteq\{1,\ldots,k\}\\
                  S\text{ has no consecutive elements}}}
 (-\beta)^{\epsilon(S)}
 a^{|S|-\epsilon(S)}
 \left(\prod_{j\in S}q_jX_j\right)
 \left(\prod_{i\in U(S)}(x_i+z)\right).
\end{equation}
In particular, no inverse power of $\beta$ is needed.
For $k\geq2$ it satisfies
\begin{equation}\label{eq:block-recurrence}
\begin{split}
 (1-\beta^2q_{k-1})B_k^q
 ={}&
 \bigl(x_k+z+\beta(q_{k-1}-q_k)
       -\beta^2(q_kx_k+zq_{k-1})\bigr)B_{k-1}^q\\
 &+(1-\beta z)q_{k-1}(1-\beta^2q_k)X_k B_{k-2}^q,
\end{split}
\end{equation}
with $B_0^q=1$ and
$B_1^q=x_1+z-\beta q_1X_1$.
\end{proposition}

\begin{proof}
Multiply \eqref{eq:top-block} by $\beta^k$ and substitute
\eqref{eq:quantum-elementary}. In each product of factors
$1-Q_j$, choose the term $-Q_j$ for a set $S$ of indices.
Such a choice forces $j\in J$ and $j+1\notin J$ for each
$j\in S$, so two elements of $S$ cannot be consecutive.
Conversely these are the only restrictions on the remaining
choices of $J$.

Each forced selected position contributes $X_j$.
Each forced unselected position $j+1\leq k$ contributes $-a$.
At every other position, summing the selected and unselected
possibilities gives $X_i-a$. Thus the contribution for $S$ to
$\beta^kB_k^q$ is
\[
 (-1)^{|S|}Q^S
 \left(\prod_{j\in S}X_j\right)
 (-a)^{|S|-\epsilon(S)}
 \prod_{i\in U(S)}(X_i-a),
 \qquad Q^S=\prod_{j\in S}Q_j.
\]
Since $|U(S)|=k-2|S|+\epsilon(S)$,
$Q_j=\beta^2q_j$, and $X_i-a=\beta(x_i+z)$, this expression is
$\beta^k$ times the summand in \eqref{eq:block-matchings}.
This proves both that formula and polynomiality.

For the recurrence, introduce the auxiliary polynomials
\[
 H_0=1,\qquad H_1=x_1+z,\qquad
 H_j=(x_j+z)H_{j-1}+a q_{j-1}X_{j-1}H_{j-2}\quad(j\geq2).
\]
They enumerate tilings of $\{1,\ldots,j\}$ by single positions
of weight $x_i+z$ and adjacent pairs $\{i,i+1\}$ of weight
$a q_iX_i$. Separating the terms of
\eqref{eq:block-matchings} according to whether $k\in S$ gives
\begin{equation}\label{eq:block-auxiliary}
 B_j^q=H_j-\beta q_jX_jH_{j-1}\qquad(j\geq1).
\end{equation}
Set $T_j=(1-\beta^2q_j)X_j$.
The recurrence for $H_j$, together with
$\beta(x_j+z-\beta q_jX_j)+a=T_j$, gives
\[
 B_k^q=(x_k+z-\beta q_kX_k)B_{k-1}^q
               +q_{k-1}X_{k-1}T_kH_{k-2}.
\]
For $k\geq3$, the same calculation at $k-1$ gives
\[
 \beta B_{k-1}^q
       =T_{k-1}H_{k-2}-aB_{k-2}^q.
\]
For $k=2$ this last identity follows directly from $B_1^q$.
Multiplying the preceding expression for $B_k^q$ by
$1-\beta^2q_{k-1}$ and using this identity yields
\[
 (1-\beta^2q_{k-1})B_k^q
   =\frac{T_k-a(1-\beta^2q_{k-1})}{\beta}\,B_{k-1}^q
       +a q_{k-1}T_kB_{k-2}^q.
\]
Expanding the numerator gives exactly
\eqref{eq:block-recurrence}; it is divisible by $\beta$.
\end{proof}

At $\beta=0$, this becomes
\[
 B_k^q=(x_k+z)B_{k-1}^q+q_{k-1}B_{k-2}^q,
\]
the elementary quantum cohomological recurrence. The additional
terms in \eqref{eq:block-recurrence} are therefore essential
K-theoretic corrections, not changes in notation.

\subsection{Passing from finite flags to polynomial identities}

We will use the equivariant quantum Chevalley formula, which is
initially an identity of classes on a finite flag variety.
The following observation explains the passage back to polynomials.

\begin{proposition}[Stable detection]\label{prop:stable-detection}
At $\beta=-1$, let $P$ be a polynomial in finitely many $x_i$ and
$q_i$, with coefficients rational in finitely many $y_a$.
If its image in the localized equivariant quantum K-theory
presentation is zero in every sufficiently large rank, then $P=0$
as a polynomial.
\end{proposition}

\begin{proof}
Suppose $P\ne0$, and choose its least nonzero total quantum degree.
In the quantum-adic expansion of its image, the coefficient of each
monomial of that degree is the ordinary equivariant K-theory class
of the corresponding coefficient polynomial $f(x;y)$.
Indeed, all smaller quantum-degree coefficient polynomials of $P$
are zero. The quantum K-theory module is free on the Schubert
classes, so each such ordinary class must vanish.

Let $x_1,\ldots,x_r$ include the variables occurring in $f$, and
let $d$ bound its degree in each of these variables.
Choose a rank $N>d+r$ containing all parameters in question.
Restricting a vanishing ordinary class to every torus fixed point
gives
\[
 f(\overline y_{\sigma(1)},\ldots,
   \overline y_{\sigma(r)};y)=0
 \qquad(\sigma\in S_N).
\]
Work over the rational function field in the equivariant parameters.
The $N$ numbers $\overline y_a$ are distinct. After fixing distinct
values for $r-1$ variables, the remaining variable has
$N-r+1>d$ possible values. Polynomial interpolation says that
every coefficient in that variable vanishes. Repeating this
argument in the other variables proves $f=0$, a contradiction.
\end{proof}

We use here the presentation and Schubert representatives of
\cite{maeno2024presentation2}. Localizing the equivariant coefficient ring does not
affect the conclusion, since clearing its denominators is injective.

\subsection{The line Chevalley formula}
\label{sec:line-chevalley}

The proof of Theorem~\ref{prop:top} is an induction on $k$ driven by
multiplication by the single line class
\[
 T_k=(1-\beta^2q_k)X_k=(1-\beta^2q_k)(1+\beta x_k),
\]
exactly as the proof of the quantum Schubert Pieri formula is driven
by multiplication by $x_k-z$. The following lemma is the
$K$-theoretic single-variable Monk formula.

\begin{definition}[Line paths]\label{def:line-path}
Let $k\geq1$ and $u\in S_\infty$, and put $M=\max\{n(u),k\}+1$ as in
Proposition~\ref{prop:rank}. A \emph{line path at $k$ from $u$} is a
directed path $\pi$ in the quantum Bruhat graph with labels
\[
 (a_1,k),\,(a_2,k),\,\ldots,(a_s,k),\qquad
 (k,b_1),\,(k,b_2),\,\ldots,(k,b_t),
\]
where $s,t\geq0$, $k>a_1>a_2>\cdots>a_s\geq1$, and
$M\geq b_1>b_2>\cdots>b_t>k$. Write $s(\pi)$ and $t(\pi)$ for the two
lengths, $v(\pi)$ for the permutation reached after the first
$s(\pi)$ edges, $\operatorname{end}(\pi)$ for the final permutation,
and $D(\pi)$ for the quantum degree, defined edgewise as in
Definition~\ref{def:chain}. The empty path is allowed, with
$v(\pi)=\operatorname{end}(\pi)=u$.
\end{definition}

\begin{lemma}[Line Chevalley formula]\label{lem:line-chevalley}
For every $k\geq1$ and $u\in S_\infty$,
\begin{equation}\label{eq:line-chevalley}
 T_k\,\G_u^q(x;y)
 =\sum_{\pi}
   (-\beta)^{s(\pi)}\beta^{t(\pi)}q^{D(\pi)}\,
   \frac{1}{1+\beta y_{v(\pi)(k)}}\,
   \G_{\operatorname{end}(\pi)}^q(x;y),
\end{equation}
the sum running over all line paths at $k$ from $u$.
\end{lemma}

\begin{proof}
Set $\beta=-1$ first and write $Y_a=(1-y_a)^{-1}$, so that the
claimed coefficient is $(-1)^{t(\pi)}q^{D(\pi)}Y_{v(\pi)(k)}$.
Work in $QK_T(Fl_N)$ for $N$ large, with the presentation and
Schubert representatives of \cite{maeno2024presentation2}: the residue class of $T_k$ is
the line bundle class $[\mathcal O(-\epsilon_k)]$ with
$\epsilon_k=\varpi_k-\varpi_{k-1}$, $\G_w^q$ represents the opposite
Schubert class $[\mathcal O^w]$, and the equivariant parameters are
identified by $1-y_a=\mathbf e^{-\epsilon_a}$.

Kato's isomorphism $\Phi$ onto the equivariant $K$-group of the
semi-infinite flag manifold sends $\mathbf e^{\mu}[\mathcal O^w]Q^\xi$
to $\mathbf e^{-\mu}[\mathcal O_{\mathbf Q}(wt_\xi)]$, and by
\cite[Section~3.1]{maeno2024presentation2}, quoting \cite[Proposition~5.2]{maeno2023presentation1}, it
intertwines quantum multiplication by $[\mathcal O(-\epsilon_k)]$ with
$(1-\mathbf t_{k-1})\,[\mathcal O_{\mathbf Q}(w_0\epsilon_k)]\otimes(-)$,
where $\mathbf t_{k-1}$ corresponds to $Q_{k-1}$ and $\mathbf t_0=0$.
The tensor product with $[\mathcal O_{\mathbf Q}(w_0\epsilon_k)]$ is
computed by the general Chevalley formula of Lenart--Naito--Sagaki
\cite[Theorem~33]{lenart2024chevalley}, in the form recalled in
\cite[Section~5.1]{maeno2024presentation2}, for the weight $\lambda=-\epsilon_k$ and
any reduced $\lambda$-chain. We use the chain of
\cite[Section~5.2]{maeno2024presentation2}: the alcove path from $A_\circ$ to
$A_{\epsilon_k}$ crosses exactly the hyperplanes
\[
 H_{\epsilon_a-\epsilon_k,\,0}\ (a=k-1,\ldots,1),\qquad
 H_{\epsilon_k-\epsilon_b,\,1}\ (b=N,\ldots,k+1),
\]
in this order. The first group is crossed in the negative direction,
so the corresponding roots of the chain are the positive roots
$\epsilon_a-\epsilon_k$; the second group is crossed in the positive
direction, so the corresponding roots are the negative roots
$-(\epsilon_k-\epsilon_b)$. An admissible subset is therefore a line
path $\pi$, and the statistic $n(A)$ counting negative roots is
$t(\pi)$. The affine reflections of the first group are linear, while
those of the second group fix $\epsilon_k$, because
$\epsilon_k\in H_{\epsilon_k-\epsilon_b,1}$; hence the weight
statistic is $-v(\pi)\epsilon_k=-\epsilon_{v(\pi)(k)}$. Finally,
$\langle-\epsilon_k,\alpha_{k-1}^\vee\rangle=1$ and all other pairings
of $-\epsilon_k$ with simple coroots are nonpositive, so the auxiliary
sum over partitions in the general formula is a geometric series
contributing $(1-\mathbf t_{k-1})^{-1}$, which cancels the factor
$1-\mathbf t_{k-1}$ above. Transporting back through $\Phi$, which
negates the exponent of $\mathbf e$ and converts $t_\xi$ into
$Q^\xi$, gives
\[
 T_k\star[\mathcal O^u]
 =\sum_\pi(-1)^{t(\pi)}Q^{D(\pi)}\mathbf e^{\epsilon_{v(\pi)(k)}}
       [\mathcal O^{\operatorname{end}(\pi)}]
 =\sum_\pi(-1)^{t(\pi)}q^{D(\pi)}Y_{v(\pi)(k)}
       [\mathcal O^{\operatorname{end}(\pi)}]
\]
in $QK_T(Fl_N)$. For $k=1$ there are no left edges and this is the
first fundamental weight case of the Chevalley formula in
\cite[Section~5]{lenart2024chevalley}, after the substitution
$[\mathcal O^{s_1}]=1-\mathbf e^{-\varpi_1}T_1$.

The bound $b_1\leq M$ is proved as in Proposition~\ref{prop:rank}:
$v(\pi)$ agrees with $u$ outside $\{1,\ldots,k\}$, so if $b_1>M$ then
$v(\pi)(b_1)=b_1$, $v(\pi)(b_1-1)=b_1-1$ and $v(\pi)(k)<b_1-1$, and
$(k,b_1)$ is neither a Bruhat nor a quantum edge. Thus the right side
of \eqref{eq:line-chevalley} is independent of $N$, and
Proposition~\ref{prop:stable-detection} upgrades the identity in
$QK_T(Fl_N)$ to a polynomial identity at $\beta=-1$. The
homogenization to general $\beta$ is the scaling used in
Proposition~\ref{prop:k-one} below: with $c=-\beta$ and
$(x,y,q)\mapsto(cx,cy,\beta^2q)$ one has $T_k\mapsto T_k$,
$Y_a\mapsto(1+\beta y_a)^{-1}$, $\G^q_w\mapsto c^{\ell(w)}\G^q_w$,
and
\[
 c^{\ell(\operatorname{end}(\pi))-\ell(u)}\beta^{2|D(\pi)|}(-1)^{t(\pi)}
 =(-\beta)^{s(\pi)}\beta^{t(\pi)}
\]
by the length--degree identity
$\ell(\operatorname{end}(\pi))-\ell(u)=s(\pi)+t(\pi)-2|D(\pi)|$.
\end{proof}

We record the operator form in which the lemma is used. For a formal
combination $\Theta=\sum_v\theta_v\,\G_v^q$ whose coefficients are
free of $x$, put
\begin{equation}\label{eq:line-operator}
 \mathcal T_k\Theta
 :=\sum_v\theta_v\sum_{\pi\text{ from }v}
    (-\beta)^{s(\pi)}\beta^{t(\pi)}q^{D(\pi)}\,
    \frac{\G^q_{\operatorname{end}(\pi)}}{1+\beta y_{v(\pi)(k)}},
\end{equation}
so that Lemma~\ref{lem:line-chevalley} reads
$T_k\Theta=\mathcal T_k\Theta$ as polynomials.

\subsection{Proof of the triple formula for \texorpdfstring{$k=1$}{k=1}}

\begin{proposition}\label{prop:k-one}
Theorem~\ref{prop:top} holds for $k=1$ and arbitrary
$u\in S_\infty$.
\end{proposition}

\begin{proof}
Take $k=1$ in Lemma~\ref{lem:line-chevalley}. There are no left
edges, so $v(\pi)=u$ for every line path, the equivariant factor is
the constant $(1+\beta y_{u(1)})^{-1}$, and
\begin{equation}\label{eq:first-line-beta}
 (1-\beta^2q_1)X_1\G_u^q
   =\frac{1}{1+\beta y_{u(1)}}
       \sum_{\gamma}\beta^{s(\gamma)}
           q^{D(\gamma)}\G_{\operatorname{end}(\gamma)}^q,
\end{equation}
where $\gamma$ runs over directed paths with labels forming a
subsequence of $(1,M),(1,M-1),\ldots,(1,2)$.
These are exactly the admissible chains of
Definition~\ref{def:chain} for $k=1$: all lower indices are $1$, and
distinctness and weak decrease of the upper indices make them
strictly decreasing.
There is at most one such path to a given endpoint.
In fact, a nonempty product
$t_{1b_1}\cdots t_{1b_s}$ with distinct $b_i$ is a single
cycle on $\{1,b_1,\ldots,b_s\}$.
The cycle of $u^{-1}w$ determines the ordered list of these
indices. It also shows that a nonempty path cannot have endpoint
$u$, and that the value $u(1)$ has moved out of position $1$.
Consequently a nonempty path has $m=1$,
$\Fix=\Left=\varnothing$, and $\Out=\{u(1)\}$.

Finally, $\beta B_1^q=(1-\beta^2q_1)X_1-(1-\beta z)$.
Subtract the second term from \eqref{eq:first-line-beta}.
The empty-path coefficient, after division by $\beta$, is
\[
 \frac{(1+\beta y_{u(1)})^{-1}-(1-\beta z)}{\beta}
       =z-\frac{y_{u(1)}}{1+\beta y_{u(1)}}.
\]
Every nonempty path has coefficient
\[
 q^{D(\gamma)}
       \frac{\beta^{s(\gamma)-1}}{1+\beta y_{u(1)}}.
\]
These are precisely \eqref{eq:coefficient} in the two cases.
All exponents are nonnegative, so clearing equivariant
denominators extends the result to $\beta=0$ as well.
\end{proof}

For later use we make the scaling in the proof of
Lemma~\ref{lem:line-chevalley} explicit. Put $c=-\beta$ and replace
$(x,y,z,q)$ at $\beta=-1$ by $(cx,cy,cz,\beta^2q)$. The defining
divided differences and the quantization map give
\begin{equation}\label{eq:scaling}
 \G_w^{\,\beta^2q}(cx;cy)\big|_{\beta=-1}
       =c^{\ellength(w)}\G_w^q(x;y),
\end{equation}
where on the left the quantum arguments are $\beta^2q$ and the
Grothendieck parameter is fixed at $-1$. The scaling takes place with
$\beta$ inverted; whenever both sides of a resulting identity have no
negative powers of $\beta$, it holds before that inversion as well.

\subsection{What the published marked-chain formula proves for all \texorpdfstring{$k$}{k}}

There is also an unconditional all-$k$ calculation, but it is
non-equivariant and initially indexed by paths rather than endpoints.
For a nonempty admissible chain $\gamma$, let
$r(\gamma)$ be the number of distinct lower indices $a_i$.
Let $F(\gamma)$ be the following set of edge indices:
\begin{itemize}
 \item every $i$ in the maximal initial segment with
       $b_1=\cdots=b_i$ and $a_1>\cdots>a_i$;
 \item every $i<s(\gamma)$ for which
       $(a_i,b_i)\succ(a_{i+1},b_{i+1})$.
\end{itemize}
Put $h(\gamma)=|F(\gamma)|$.
These are exactly the indices forced to be marked in the
Naito--Sagaki marking rule; in particular $h(\gamma)\geq1$.

\begin{proposition}\label{prop:marked-block}
For every $k\geq1$ and $u\in S_\infty$,
\begin{equation}\label{eq:marked-block}
\begin{split}
 B_k^q(x;z)\G_u^q(x;0)
  ={}&z^k\G_u^q(x;0)\\
   &+\sum_{\substack{\gamma\text{ admissible}\\s(\gamma)>0}}
       q^{D(\gamma)}
       \beta^{s(\gamma)-r(\gamma)}
       z^{k-r(\gamma)}
       (1-\beta z)^{h(\gamma)-1}
       \G_{\operatorname{end}(\gamma)}^q(x;0).
\end{split}
\end{equation}
This statement does not require uniqueness of a path to its endpoint.
\end{proposition}

\begin{proof}
For $1\leq p\leq k$, let
$c[k,p]=(k-p+1,k-p+2,\ldots,k+1)$, in cycle notation.
The single Grothendieck polynomial of this cycle is
\[
 \G_{c[k,p]}(x;0)
   =\sum_{j=p}^k
       \binom{j-1}{p-1}\beta^{j-p}e_j(x_1,\ldots,x_k).
\]
The binomial theorem, coefficient by coefficient in $e_j$, gives
\[
 \prod_{i=1}^k(x_i+z)
  =z^k+\sum_{p=1}^k
       z^{k-p}(1-\beta z)^{p-1}\G_{c[k,p]}(x;0).
\]
Indeed the coefficient of $e_j$ on the right is
$z^{k-j}\bigl((1-\beta z)+\beta z\bigr)^{j-1}=z^{k-j}$.
Applying the \emph{linear} map $\Q$ gives
\begin{equation}\label{eq:block-cyclic}
 B_k^q=z^k+\sum_{p=1}^k
       z^{k-p}(1-\beta z)^{p-1}\G_{c[k,p]}^q(x;0).
\end{equation}
No multiplicativity of $\Q$ is used.

Naito--Sagaki's Pieri theorem \cite[Theorem~2.13]{naito2024pieri},
homogenized from $\beta=-1$, says that the contribution of
$\gamma$ to
$\G_u^q\G_{c[k,p]}^q$ is
\[
 q^{D(\gamma)}\beta^{s(\gamma)-p}
       \#\operatorname{Mark}_p(\gamma).
\]
A mark may occur only at a first occurrence of a lower index.
All indices in $F(\gamma)$ must be marked; the remaining first
occurrences can be chosen freely. The admissibility conditions
ensure that forced indices are first occurrences.
Thus, as also recorded in \cite[Remark~2.14]{naito2024pieri},
\[
 \#\operatorname{Mark}_p(\gamma)
       =\binom{r(\gamma)-h(\gamma)}{p-h(\gamma)}.
\]
Here $1\leq h(\gamma)\leq r(\gamma)\leq k$.

Fix a nonempty chain, and abbreviate its statistics by $s,r,h,D$.
Its coefficient after multiplying \eqref{eq:block-cyclic} by
$\G_u^q$ is therefore
\begin{align*}
 &q^D\sum_{p=h}^r
       \binom{r-h}{p-h}
       \beta^{s-p}z^{k-p}(1-\beta z)^{p-1}\\
 &\qquad=
 q^D\beta^{s-r}z^{k-r}(1-\beta z)^{h-1}
 \sum_{p=h}^r\binom{r-h}{p-h}
       (\beta z)^{r-p}(1-\beta z)^{p-h}\\
 &\qquad=
 q^D\beta^{s-r}z^{k-r}(1-\beta z)^{h-1}.
\end{align*}
The last sum is $1$ by the binomial theorem.
The separate term $z^k\G_u^q$ supplies the empty path.
This proves \eqref{eq:marked-block}.
Homogenization is justified by the same scaling and
length--degree identity used in Proposition~\ref{prop:k-one};
the exponents $s-p$ are nonnegative.
\end{proof}

With Corollary~\ref{cor:rotation} in hand, the path-indexed sum
\eqref{eq:marked-block} is in fact indexed by endpoints, and its
exponents can be read off from the endpoint. Write $\gamma_w$ for the
unique admissible $k$-chain from $u$ to $w\in E(u,k)$.

\begin{lemma}[Forced marks count left moves]\label{lem:h-left}
Let $w\in E(u,k)$, $w\neq u$. Then
\[
 r(\gamma_w)=m(u,w;k)
 \qquad\text{and}\qquad
 h(\gamma_w)=|\Left(u,w;k)|+1 .
\]
\end{lemma}

\begin{proof}
The first equality is Corollary~\ref{cor:rotation}(2). For the
second, note first that $(a_i,b_i)\succ(a_{i+1},b_{i+1})$ means
$b_i=b_{i+1}$ and $a_i>a_{i+1}$; call such an index $i$ a
\emph{descent} of $\gamma_w$. If $i_0$ is the last index of the
maximal initial segment with equal upper indices and decreasing lower
indices, then $1,\ldots,i_0-1$ are descents and $i_0$ is not, so
$F(\gamma_w)$ is the disjoint union of $\{i_0\}$ and the set of
descents, and $h(\gamma_w)-1$ is the number of descents. We show that
the descents are in bijection with $\Left(u,w;k)$.

By Corollary~\ref{cor:rotation} we may rotate the values and assume
that $\gamma_w$ is a saturated chain in the $k$-Bruhat order
satisfying (P0) and (P1); the sets $\Fix,\Left,\Out$ and the
descents only involve positions and are unchanged. Consider a
maximal run $(a_1,b),\ldots,(a_p,b)$ of labels with the same upper
index $b$, and let $\alpha_j$ be the value at position $a_j$ just
before the label $(a_j,b)$ is applied, $\alpha_0$ the value at $b$
before the run. Each label is a Bruhat cover, so
$\alpha_0>\alpha_1>\cdots>\alpha_p$, and the net effect of the run is
that $\alpha_0$ moves to $a_1$, $\alpha_j$ moves from $a_j$ to
$a_{j+1}$ for $1\leq j<p$, and $\alpha_p$ moves from $a_p$ to $b$.
Position $b$ is touched by no other label, so $\alpha_p$ stays at
$b$ for ever, and every value entering the window from outside does
so at the start of a run. In particular a value moves within the
window only in steps $a_j\to a_{j+1}$, which is a move to the left
exactly at a descent.

Let $i$ be a descent, so that the labels $i$ and $i+1$ are consecutive
labels $(a_j,b),(a_{j+1},b)$ of a run with $a_{j+1}<a_j$, and the
value $\alpha_j$ moves from $a_j$ to $a_{j+1}$. By (P0) the lower index
$a_j$ is a first occurrence, so position $a_j$ was untouched and
$\alpha_j=u(a_j)$ is a window value making its first move. By the
prohibition lemma of Lenart and Sottile
\cite[Lemma~2.9]{lenart2007pieri}, a chain satisfying (P0) and (P1)
contains no subsequence $(j',m),\ldots,(i',m),\ldots,(i',l)$ with
$i'<j'\leq k<l<m$; applied with $(j',m)=(a_j,b)$ and
$(i',m)=(a_{j+1},b)$, this says that position $a_{j+1}$ is never
touched again (a later label $(a_{j+1},l)$ has $l<b$ by (P1) and
distinctness). Hence $u(a_j)$ ends at position $a_{j+1}<a_j$ and lies
in $\Left(u,w;k)$. Distinct descents move distinct values, so this
gives an injection from descents into $\Left(u,w;k)$.

Conversely let $c=u(i)\in\Left(u,w;k)$. It moves at least once;
its first move is by a label $(i,b)$. If this label ends its run,
$c$ stays at $b$ and is in $\Out$; otherwise $c$ moves to the lower
index $a'$ of the next label. If $a'<i$ this is a descent, and $c$ is
the value attached to it. If $a'>i$, then every later move of $c$
starts from a position that has already been touched, so by (P0) it
is a step to a larger lower index or the end of a run; thus $c$ only
moves right or leaves the window, and ends in $\Out$, a
contradiction. Hence every value in $\Left(u,w;k)$ is attached to a
descent, and the correspondence is a bijection.
\end{proof}

\begin{theorem}[The non-equivariant case]\label{thm:noneq}
Theorem~\ref{prop:top} holds at $y=0$: for every $k\geq1$ and
$u\in S_\infty$,
\[
 B_k^q(x;z)\,\G_u^q(x;0)
 =\sum_{w\in E(u,k)}
   q^{D(u,w;k)}\beta^{s(u,w;k)-m(u,w;k)}
   z^{|\Fix(u,w;k)|}(1-\beta z)^{|\Left(u,w;k)|}\,
   \G_w^q(x;0).
\]
\end{theorem}

\begin{proof}
By Corollary~\ref{cor:rotation}(1) the sum in
\eqref{eq:marked-block} has exactly one term for each
$w\in E(u,k)\setminus\{u\}$, namely $\gamma=\gamma_w$, and the term
$z^k\G_u^q$ is the term $w=u$ with $\Fix=\{u(1),\ldots,u(k)\}$.
By Lemma~\ref{lem:h-left}, $k-r(\gamma_w)=k-m=|\Fix|$,
$s(\gamma_w)-r(\gamma_w)=s-m$ and $h(\gamma_w)-1=|\Left|$. This is
\eqref{eq:coefficient} at $y=0$.
\end{proof}

Thus the only part of Theorem~\ref{prop:top} that remains open is
the dependence on the equivariant parameters $y$. At $z=1/\beta$ this
dependence is settled in Theorem~\ref{thm:a-zero} below.

\subsection{Reduction of level \texorpdfstring{$k$}{k} to levels \texorpdfstring{$k-1$}{k-1} and \texorpdfstring{$k-2$}{k-2}}
\label{sec:reduction}

We now set up the induction on $k$. Throughout, $a=1-\beta z$ and
$H_j$ denotes the tiling polynomial introduced in the proof of
Proposition~\ref{prop:block-polynomial}:
\[
 H_0=1,\qquad H_1=x_1+z,\qquad
 H_j=(x_j+z)H_{j-1}+a\,q_{j-1}X_{j-1}H_{j-2}\quad(j\geq2),
\]
and $B_j^q=H_j-\beta q_jX_jH_{j-1}$ for $j\geq1$ by
\eqref{eq:block-auxiliary}. We use the convention $q_0=0$ and
$B_{-1}^q=0$.

\begin{lemma}[Two-step block identities]\label{lem:two-step}
For every $k\geq1$,
\begin{align}
 \beta B_k^q+aB_{k-1}^q&=T_kH_{k-1},\label{eq:step-one}\\
 (1-\beta^2q_{k-1})H_{k-1}&=B_{k-1}^q+a\beta q_{k-1}B_{k-2}^q.
 \label{eq:step-two}
\end{align}
Consequently
\begin{equation}\label{eq:step-combined}
 (1-\beta^2q_{k-1})\bigl(\beta B_k^q+aB_{k-1}^q\bigr)
 =T_k\bigl(B_{k-1}^q+a\beta q_{k-1}B_{k-2}^q\bigr).
\end{equation}
\end{lemma}

\begin{proof}
Substituting $B_k^q=H_k-\beta q_kX_kH_{k-1}$,
$B_{k-1}^q=H_{k-1}-\beta q_{k-1}X_{k-1}H_{k-2}$ and the recurrence for
$H_k$,
\[
 \beta B_k^q+aB_{k-1}^q
 =\bigl(\beta(x_k+z)+a-\beta^2q_kX_k\bigr)H_{k-1}
 =(1-\beta^2q_k)X_kH_{k-1},
\]
since $\beta(x_k+z)+a=1+\beta x_k=X_k$; the two terms involving
$a\beta q_{k-1}X_{k-1}H_{k-2}$ cancel. This is \eqref{eq:step-one}.
For \eqref{eq:step-two} with $k\geq3$, expand both sides using the
recurrence for $H_{k-1}$ and the expressions of $B_{k-1}^q,B_{k-2}^q$
through $H$: the right side is
\[
 H_{k-1}+\beta q_{k-1}(a-X_{k-1})H_{k-2}
   -a\beta^2q_{k-1}q_{k-2}X_{k-2}H_{k-3},
\]
and $\beta(a-X_{k-1})=-\beta^2(x_{k-1}+z)$, so this equals
$H_{k-1}-\beta^2q_{k-1}\bigl((x_{k-1}+z)H_{k-2}+aq_{k-2}X_{k-2}H_{k-3}\bigr)
=(1-\beta^2q_{k-1})H_{k-1}$. For $k=2$ the same computation applies
without the last term, and for $k=1$ both sides equal $1$.
Multiplying \eqref{eq:step-one} by $1-\beta^2q_{k-1}$ and
substituting \eqref{eq:step-two} gives \eqref{eq:step-combined}.
\end{proof}

At $\beta=0$, \eqref{eq:step-combined} reduces to
$B_k=(x_k+z)B_{k-1}+q_{k-1}B_{k-2}$, the recurrence used for the
quantum Schubert Pieri formula. At $\beta=-1$ it reads
$(1-q_{k-1})\bigl(AB_{k-1}-B_k\bigr)=T_k\bigl(B_{k-1}-Aq_{k-1}B_{k-2}\bigr)$
with $A=1+z$.

\paragraph{The target identity.}
For $j\geq0$ and $u\in S_\infty$ write
\[
 P_j(u):=\sum_{w\in E(u,j)}C_{u,w}^{(j)}\,\G_w^q,
 \qquad P_0(u)=\G_u^q,\qquad P_{-1}(u)=0,
\]
for the right side of \eqref{eq:main} at level $j$, with
$C^{(j)}_{u,w}$ given by \eqref{eq:coefficient}. The endpoint
data $(s,D)$ entering $C^{(j)}_{u,w}$ are well defined by
Corollary~\ref{cor:rotation}.

\begin{proposition}[Reduction]\label{prop:reduction}
Let $k\geq1$ and suppose Theorem~\ref{prop:top} holds at levels
$k-1$ and $k-2$ for the permutation $u$, that is,
$B_{k-1}^q\G_u^q=P_{k-1}(u)$ and $B_{k-2}^q\G_u^q=P_{k-2}(u)$.
Then Theorem~\ref{prop:top} holds at level $k$ for $u$ if and only if
\begin{equation}\label{eq:junction}
 \beta(1-\beta^2q_{k-1})\,P_k(u)
 =\mathcal T_k\bigl(P_{k-1}(u)+a\beta q_{k-1}P_{k-2}(u)\bigr)
   -a(1-\beta^2q_{k-1})\,P_{k-1}(u),
\end{equation}
with $\mathcal T_k$ the line operator \eqref{eq:line-operator}.
\end{proposition}

\begin{proof}
By \eqref{eq:step-combined} and the hypotheses,
\[
 (1-\beta^2q_{k-1})\bigl(\beta B_k^q\G_u^q+aP_{k-1}(u)\bigr)
 =T_k\bigl(P_{k-1}(u)+a\beta q_{k-1}P_{k-2}(u)\bigr),
\]
and the right side equals
$\mathcal T_k\bigl(P_{k-1}(u)+a\beta q_{k-1}P_{k-2}(u)\bigr)$ by
Lemma~\ref{lem:line-chevalley}. Hence
$\beta(1-\beta^2q_{k-1})B_k^q\G_u^q$ equals the right side of
\eqref{eq:junction}. Since the $\G_w^q$ are linearly independent over
the coefficient field and $\beta(1-\beta^2q_{k-1})$ is a nonzero
element of the integral domain of coefficients,
$B_k^q\G_u^q=P_k(u)$ holds if and only if \eqref{eq:junction} does.
\end{proof}

For $k=1$ the identity \eqref{eq:junction} is
$\beta P_1(u)=\mathcal T_1\G_u^q-a\G_u^q$, which is exactly the
computation in the proof of Proposition~\ref{prop:k-one}. Thus
Theorem~\ref{prop:top} follows, by induction on $k$, from
\eqref{eq:junction} for all $k\geq2$ and all $u$. We call
\eqref{eq:junction} the \emph{junction identity}; it is the
$K$-theoretic counterpart of the recurrence
$\mathcal E_{k,k}=(x_k-z_1)\mathcal E_{k-1,k-1}+q_{k-1}\mathcal E_{k-2,k-2}$ for quantum factorial elementary polynomials
expanded by the single-variable Monk formula in the quantum Schubert
case.

\paragraph{Coefficient form.}
Extracting the coefficient of $\G_w^q$ in \eqref{eq:junction}, and
writing $\pi\colon v\to w$ for a line path at $k$ from $v$ with
endpoint $w$, the junction identity reads
\begin{equation}\label{eq:junction-coeff}
\begin{split}
 \beta(1-\beta^2q_{k-1})\,C^{(k)}_{u,w}
 ={}&\sum_{v}\ \sum_{\pi\colon v\to w}
     \frac{(-\beta)^{s(\pi)}\beta^{t(\pi)}q^{D(\pi)}}
          {1+\beta y_{v(\pi)(k)}}
     \Bigl(C^{(k-1)}_{u,v}+a\beta q_{k-1}C^{(k-2)}_{u,v}\Bigr)\\
   &-a(1-\beta^2q_{k-1})\,C^{(k-1)}_{u,w},
\end{split}
\end{equation}
where $C^{(j)}_{u,v}=0$ for $v\notin E(u,j)$ and $C^{(-1)}=0$.
A line path $\pi\colon v\to w$ is determined by its label sequence,
and $v=w\,t_{kb_t}\cdots t_{kb_1}t_{a_sk}\cdots t_{a_1k}$; the path
must be directed from $v$ in the quantum Bruhat graph. The term
$\pi=\varnothing$, $v=w$, is the diagonal term, with equivariant
factor $(1+\beta y_{w(k)})^{-1}$; the paths with $s(\pi)=0$ are the
up-neighbours $v=wt_{kb_t}\cdots t_{kb_1}$, and the paths with
$t(\pi)=0$ are the down-neighbours $v=wt_{a_sk}\cdots t_{a_1k}$. In
the quantum Schubert case only single edges occur and
\eqref{eq:junction-coeff} becomes the four-term expansion of that
proof; here a path may combine several down edges followed by several
up edges, and each term carries a $K$-theoretic sign and an
equivariant denominator determined by the intermediate permutation
$v(\pi)$.

\begin{remark}[Well-definedness]\label{rem:well-defined}
The identity \eqref{eq:junction-coeff} is to be proved for each
endpoint $w$. By Corollary~\ref{cor:rotation} the data $(s,D)$
entering $C^{(k)}_{u,w}$ are already known to be well defined, and
$D$ is given by \eqref{eq:D-closed}; so the left side is unambiguous,
and a proof of \eqref{eq:junction-coeff} need not establish
path-independence separately. It does, however, have to show that the
right side collapses to the single monomial $q^{D(u,w;k)}$ times the
expected endpoint factors, and that it vanishes for $w\notin E(u,k)$.
In the quantum Schubert proof the analogous path-independence was
proved through the cycle structure of $u^{-1}w$; here it comes from
the rotation to the classical $k$-Bruhat order.
\end{remark}

\paragraph{The diagonal case.}
Suppose $u(k)=w(k)$. Position $k$ is touched by a $k$-chain only
through labels $(k,b)$, and by Corollary~\ref{cor:rotation}(2) a
touched position never regains its original value. Hence the
admissible $k$-chains from $u$ to $w$ have no label with lower index
$k$; they are exactly the admissible $(k-1)$-chains from $u$ to $w$
all of whose labels have second entry at least $k+1$,
and for these chains the conditions of Definition~\ref{def:chain} at
levels $k$ and $k-1$ coincide. Corollary~\ref{cor:rotation}(2) at
level $k-1$, where $k$ is an outside position, shows likewise
that the $(k-1)$-chains to $w$ with some label $(a,k)$ have
$w(k)\neq u(k)$. Therefore
\[
 \Fix(u,w;k)=\Fix(u,w;k-1)\sqcup\{u(k)\},\qquad
 \Left,\ \Out,\ s,\ D\ \text{unchanged},
\]
and
\begin{equation}\label{eq:diag-step}
 C^{(k)}_{u,w}=\Bigl(z-\frac{y_{u(k)}}{1+\beta y_{u(k)}}\Bigr)
               C^{(k-1)}_{u,w}
 =\frac{(1+\beta y_{u(k)})^{-1}-a}{\beta}\,C^{(k-1)}_{u,w}.
\end{equation}
Inserting \eqref{eq:diag-step} into \eqref{eq:junction-coeff}, the
diagonal term $\pi=\varnothing$ together with the last term produce
exactly $\beta C^{(k)}_{u,w}$ plus
$\frac{a\beta q_{k-1}}{1+\beta y_{u(k)}}C^{(k-2)}_{u,w}$,
and what remains to be shown in the case $u(k)=w(k)$ is the
cancellation
\begin{equation}\label{eq:diag-cancel}
 \sum_{v}\ \sum_{\substack{\pi\colon v\to w\\ \pi\neq\varnothing}}
     \frac{(-\beta)^{s(\pi)}\beta^{t(\pi)}q^{D(\pi)}}
          {1+\beta y_{v(\pi)(k)}}
     \Bigl(C^{(k-1)}_{u,v}+a\beta q_{k-1}C^{(k-2)}_{u,v}\Bigr)
 +\frac{\beta q_{k-1}}{1+\beta y_{u(k)}}
     \Bigl(\beta C^{(k-1)}_{u,w}+aC^{(k-2)}_{u,w}\Bigr)=0,
\end{equation}
where we used $\beta(z+\overline y_{u(k)})+a=(1+\beta y_{u(k)})^{-1}$
to combine the two multiples of $q_{k-1}C^{(k-1)}_{u,w}$.
This is the $K$-theoretic form of the reflection--drop cancellation
of the quantum Schubert proof: there the nonempty paths are single
edges $t_{ak}$, $t_{kb}$ and the drop term is $q_{k-1}c_{k-2}(u,w)$.
Note the combination $\beta C^{(k-1)}_{u,w}+aC^{(k-2)}_{u,w}$, the
coefficient analogue of $\beta B_{k-1}^q+aB_{k-2}^q=T_{k-1}H_{k-2}$.

\subsection{The case \texorpdfstring{$z=1/\beta$}{z=1/beta}: the Chevalley formula for \texorpdfstring{$\varpi_k$}{varpi\_k}}
\label{sec:a-zero}

At $z=1/\beta$, that is $a=0$, the junction identity loses its
level-$(k-2)$ part and its last term, and we can prove it completely.
This gives the full equivariant statement of Theorem~\ref{prop:top}
on the hyperplane $a=0$, for all $k$ and $u$. Since
$x_i+1/\beta=X_i/\beta$, the block is
$B_k^q|_{z=1/\beta}=\beta^{-k}\Q(X_1\cdots X_k)=\beta^{-k}F_k^k
=\beta^{-k}(1-\beta^2q_k)X_1\cdots X_k$, and the coefficient
\eqref{eq:coefficient} specializes by
$(1-\beta z)^{|\Left|}=0^{|\Left|}$ and
$z+\overline y_c=(1+\beta y_c)^{-1}/\beta$ to
\[
 C^{(k)}_{u,w}\big|_{z=1/\beta}
 =\begin{cases}
   q^{D}\beta^{s-k}\prod_{i=1}^k(1+\beta y_{u(i)})^{-1}
    & \Left(u,w;k)=\varnothing,\\
   0&\text{otherwise},
  \end{cases}
\]
using $m+|\Fix|=k$ and $\Fix\sqcup\Out=\{u(1),\ldots,u(k)\}$ when
$\Left=\varnothing$. Thus the statement to be proved is the
following.

\begin{theorem}[Theorem~\ref{prop:top} at $z=1/\beta$]\label{thm:a-zero}
For every $k\geq1$ and $u\in S_\infty$,
\begin{equation}\label{eq:a-zero}
 (1-\beta^2q_k)\,X_1X_2\cdots X_k\,\G^q_u(x;y)
 =\prod_{i=1}^k\frac1{1+\beta y_{u(i)}}
  \sum_{\substack{w\in E(u,k)\\ \Left(u,w;k)=\varnothing}}
   \beta^{s(u,w;k)}q^{D(u,w;k)}\,\G^q_w(x;y).
\end{equation}
\end{theorem}

The polynomial $(1-\beta^2q_k)X_1\cdots X_k=F^k_k$ represents, in
$QK_T(Fl_N)$ with the identifications used in the proof of
Lemma~\ref{lem:line-chevalley}, the class of the line bundle
$\mathcal O(-\varpi_k)$: by \cite[Section~3.1]{maeno2024presentation2},
$[\mathcal O(-\varpi_k)]$ is the quantum product of
$[\mathcal O(-\epsilon_1)],\frac1{1-Q_1}[\mathcal O(-\epsilon_2)],
\ldots,\frac1{1-Q_{k-1}}[\mathcal O(-\epsilon_k)]$, and
$\frac{1}{1-Q_{j-1}}[\mathcal O(-\epsilon_j)]$ is represented by
$\frac{1-\beta^2q_j}{1-\beta^2q_{j-1}}X_j$. So \eqref{eq:a-zero} is a
chain-indexed form of the equivariant quantum $K$-Chevalley formula
for $\mathcal O(-\varpi_k)$, with the equivariant weight
$\prod_{i\leq k}(1+\beta y_{u(i)})^{-1}$ independent of the endpoint.
We prove it from the single-variable formula
(Lemma~\ref{lem:line-chevalley}) by the junction identity, and we
record the combinatorial lemmas separately, since they are also the
ingredients of the general case.

\paragraph{Left-free chains.}
Call an admissible $k$-chain, or its endpoint, \emph{Left-free} if
$\Left(u,w;k)=\varnothing$. Recall that $\prec$ is the total order on
labels with $(a,b)\prec(c,d)$ if $b>d$, or $b=d$ and $a<c$.

\begin{lemma}[Left-free chains]\label{lem:left-free}
An admissible $k$-chain is Left-free if and only if its label sequence
is strictly increasing for $\prec$. Conversely, every directed path in
the quantum Bruhat graph from $u$ whose labels $(a_i,b_i)$ satisfy
$a_i\leq k<b_i$ and are strictly increasing for $\prec$ is an
admissible $k$-chain, and it is Left-free.
\end{lemma}

\begin{proof}
For an admissible chain the upper indices are weakly decreasing, so
the sequence is strictly $\prec$-increasing exactly when there is no
index $i$ with $b_i=b_{i+1}$ and $a_i>a_{i+1}$, that is, no descent in
the sense of the proof of Lemma~\ref{lem:h-left}. That proof shows
that the descents are in bijection with $\Left(u,w;k)$. For the
converse, a strictly $\prec$-increasing sequence has distinct labels
and weakly decreasing upper indices, and condition~4 of
Definition~\ref{def:chain} holds for every index, repeated lower index
or not. Hence the path is an admissible chain, and it is Left-free by
the first part.
\end{proof}

Left-free chains are therefore the paths of the quantum alcove model
for the reflection order $\prec$ on the roots $\epsilon_a-\epsilon_b$,
$a\leq k<b$, which is consistent with the Chevalley interpretation of
\eqref{eq:a-zero}.

\paragraph{Three lemmas on rotated paths.}
Fix $u\in S_R$ and $k$, and let $\vartheta=u(k+1)$ and
$\rho=\rho_\vartheta$ be as in Lemma~\ref{lem:rotation}. For
$v\in S_R$ and $1\leq j<R$ put
$\sigma_j(v)=\#\{i\leq j:\ v(i)\leq\vartheta\}$.

\begin{lemma}[Degree along rotated paths]\label{lem:sigma}
Let $u=v_0\to v_1\to\cdots\to v_s=w$ be a directed path in the quantum
Bruhat graph with labels $(a_i,b_i)$, $a_i\leq k<b_i$, such that every
edge of the rotated path $\rho v_0\to\cdots\to\rho v_s$ is a Bruhat
edge. Then $D_j=\sigma_j(w)-\sigma_j(u)$ for every $j$, where $D_j$ is
the number of quantum edges of the path with $a_i\leq j<b_i$.
\end{lemma}

\begin{proof}
An edge $(a,b)$ with $a\leq j<b$ applied to $v$ changes $\sigma_j$ by
$+1$ if $v(b)\leq\vartheta<v(a)$, by $-1$ if $v(a)\leq\vartheta<v(b)$,
and not at all otherwise. By Corollary~\ref{cor:cyclic}, since the
rotated edge is Bruhat, the second case does not occur, and the first
case occurs exactly when the edge is quantum. Edges with both
positions at most $j$ or both greater than $j$ do not change
$\sigma_j$ and are not counted by $D_j$. Summing along the path gives
the claim.
\end{proof}

\begin{lemma}[Commuting a right edge]\label{lem:commute}
Let $x\in S_R$ and $a<k<b'<b$. Then
$x\lessdot xt_{kb}\lessdot xt_{kb}t_{ab'}$ is a chain of Bruhat covers
if and only if $x\lessdot xt_{ab'}\lessdot xt_{ab'}t_{kb}$ is. The two
chains have the same endpoint.
\end{lemma}

\begin{proof}
The transpositions commute, so the endpoints agree. Write
$A=x(a)$, $K=x(k)$, $B'=x(b')$ and $B=x(b)$. Recall that
$y\lessdot yt_{ij}$ means $y(i)<y(j)$ and that no value at a position
strictly between $i$ and $j$ lies in the interval $(y(i),y(j))$. The
positions strictly between $k$ and $b$ other than $b'$, and the
positions strictly between $a$ and $b'$ other than $k$, carry the same
values in all four permutations involved, so only the four values
$A,K,B',B$ matter. The first chain is a chain of covers if and only if
\[
 K<B,\quad B'\notin(K,B),\quad A<B',\quad B\notin(A,B')
\]
together with the conditions on the other positions, and the second
if and only if
\[
 A<B',\quad K\notin(A,B'),\quad K<B,\quad A\notin(K,B)
\]
together with the same conditions on the other positions. Assume the
first set. If $B'<K$, then $K\notin(A,B')$ and $A<B'<K$ give
$A\notin(K,B)$. If $B'>B$, then $B\in(K,B')$ and $B\notin(A,B')$ force
$A>B>K$, so again $K\notin(A,B')$ and $A\notin(K,B)$. Conversely assume
the second set. If $K<A$, then $A\notin(K,B)$ forces $A>B$, so
$B<A<B'$, whence $B\notin(A,B')$ and $B'\notin(K,B)$. If $K>B'$, then
$B'<K$ gives $B'\notin(K,B)$ and $B>K>B'$ gives $B\notin(A,B')$.
\end{proof}

For a sequence $\gamma$ of labels $(a,b)$ with $a\leq k<b$ let
$\operatorname{sort}(\gamma)$ be the sequence obtained by moving all
labels with lower index $k$ to the end, preserving the relative order
of these labels and of the remaining labels.

\begin{lemma}[Sorting]\label{lem:sort}
Let $\gamma$ be a strictly $\prec$-increasing sequence of labels
$(a,b)$ with $a\leq k<b$, and let $x\in S_R$. Then $\gamma$ is a chain
of Bruhat covers from $x$ if and only if $\operatorname{sort}(\gamma)$
is, and the two chains have the same endpoint.
\end{lemma}

\begin{proof}
If a label $(k,b)$ precedes a label $(a,b')$ with $a<k$ in $\gamma$,
then $(k,b)\prec(a,b')$ forces $b'<b$, since $b'=b$ would require
$k<a$. The sorted sequence is obtained from $\gamma$ by repeatedly
interchanging an adjacent pair $(k,b),(a,b')$ with $a<k$ into
$(a,b'),(k,b)$; the pairs interchanged are always in their original
relative order, so $b'<b$ holds at each step. Each step preserves the
endpoint and, by Lemma~\ref{lem:commute} applied to the permutation
reached before the pair, the property of being a chain of covers.
\end{proof}

\begin{lemma}[Rotation of sorted words]\label{lem:rotate-sorted}
Let $\omega=\gamma'\,(k,b_1)\cdots(k,b_r)$ be a directed path in the
quantum Bruhat graph from $u$, where the labels of $\gamma'$ have
lower index at most $k-1$ and upper index at least $k+1$, are
pairwise distinct with weakly decreasing upper indices, and
$b_1>\cdots>b_r>k$. Then every edge of the rotated path
$\rho u\to\cdots$ is a Bruhat edge.
\end{lemma}

\begin{proof}
The labels of $\gamma'$ satisfy the hypotheses of
Lemma~\ref{lem:rotation}, so the rotated edges of $\gamma'$ are
Bruhat. Work in the rotated order, with $N=\rho(\vartheta)$ the
largest value, initially at position $k+1$. Let
$(c_1,k+1),\ldots,(c_m,k+1)$ be the labels of $\gamma'$ with upper
index $k+1$; they are its last $m$ labels. As in the proof of
Lemma~\ref{lem:rotation}, put $e_0=N$ and let $e_h$ be the value at
position $c_h$ just before $(c_h,k+1)$; then $e_0>e_1>\cdots>e_m$,
after $\gamma'$ the value at position $k+1$ is $e_m$, and $e_{h-1}$
sits at position $c_h$. Position $k$ is touched by no label of
$\gamma'$; let $x$ be its value. Since $c_h<k<k+1$, the Bruhat
condition for $(c_h,k+1)$ gives $x\notin(e_h,e_{h-1})$ for each $h$,
and $x$ differs from all $e_h$. Hence $x\notin(e_m,N)$, and as $x<N$
we get $x<e_m$.

Let $y_0=x$ and let $y_i$ be the value at position $k$ after the edge
$(k,b_i)$. We show by induction on $i$ that $(k,b_i)$ is Bruhat and
that $e_m>y_i$ whenever $b_i\neq k+1$; the induction hypothesis is
$e_m>y_{i-1}$, true for $i=1$. Position $k+1$ still carries $e_m$ when
$(k,b_i)$ is applied, because $b_1,\ldots,b_{i-1}>b_i\geq k+1$. If
$(k,b_i)$ were quantum, every value at a position strictly between
$k$ and $b_i$ would lie in $(v(b_i),y_{i-1})$, and $v(b_i)<y_{i-1}$;
for $b_i>k+1$ this puts $e_m$ below $y_{i-1}$, and for $b_i=k+1$ it
says $e_m=v(b_i)<y_{i-1}$; both contradict $e_m>y_{i-1}$. So the edge
is Bruhat, $y_i=v(b_i)>y_{i-1}$, and if $b_i>k+1$ the Bruhat condition
gives $e_m\notin(y_{i-1},y_i)$, whence $e_m>y_i$. If $b_i=k+1$ it is
the last right edge.
\end{proof}

\begin{corollary}[Sorting bijection]\label{cor:sort-bijection}
Let $w\in S_R$. The map $\gamma\mapsto\operatorname{sort}(\gamma)$ is
a bijection from the Left-free admissible $k$-chains from $u$ to $w$
onto the directed paths from $u$ to $w$ of the form
$\omega=\gamma'\,(k,b_1)\cdots(k,b_r)$, where $\gamma'$ is a Left-free
admissible $(k-1)$-chain from $u$ none of whose labels has upper index
$k$, and $b_1>\cdots>b_r>k$. It preserves the number of edges and the
quantum degree $D$.
\end{corollary}

\begin{proof}
Let $\gamma$ be a Left-free $k$-chain from $u$ to $w$. By
Lemma~\ref{lem:rotation} the rotated path is a chain of Bruhat covers
from $\rho u$ to $\rho w$; by Lemma~\ref{lem:left-free} its labels are
strictly $\prec$-increasing, so by Lemma~\ref{lem:sort}
$\operatorname{sort}(\gamma)$ is a chain of Bruhat covers from
$\rho u$ to $\rho w$, and by Corollary~\ref{cor:cyclic} it is a
directed path from $u$ to $w$ in the quantum Bruhat graph. Its labels
with lower index less than $k$ form a strictly $\prec$-increasing
directed path $\gamma'$ from $u$ with upper indices greater than $k$,
hence a Left-free $(k-1)$-chain by Lemma~\ref{lem:left-free}, and its
labels $(k,b)$ appear in $\prec$-order, that is with decreasing $b$.
Both $\gamma$ and $\operatorname{sort}(\gamma)$ have all rotated
edges Bruhat, so by Lemma~\ref{lem:sigma} they have the same $D$, and
they have the same number of edges.

Conversely let $\omega$ be a path of the stated form. By
Lemma~\ref{lem:rotate-sorted} its rotated edges are Bruhat. Let
$\gamma$ be the $\prec$-increasing rearrangement of its labels, so
that $\operatorname{sort}(\gamma)=\omega$; the labels are distinct
since $\gamma'$ has distinct labels and the $(k,b_i)$ are distinct
from them. By Lemma~\ref{lem:sort} $\gamma$ is a chain of Bruhat
covers from $\rho u$ to $\rho w$, hence a directed path from $u$ to
$w$, and by Lemma~\ref{lem:left-free} a Left-free $k$-chain. The two
constructions are inverse to each other.
\end{proof}

\paragraph{Proof of Theorem~\ref{thm:a-zero}.}
We induct on $k$, reading the right side of \eqref{eq:a-zero} as
$\G^q_u$ for $k=0$, where the statement is trivial. Let $k\geq1$ and
assume \eqref{eq:a-zero} at level $k-1$. Since
$(1-\beta^2q_{k-1})F^k_k=T_kF^{k-1}_{k-1}$, the induction hypothesis
and Lemma~\ref{lem:line-chevalley} give
\[
 (1-\beta^2q_{k-1})F^k_k\,\G^q_u
 =\prod_{i=1}^{k-1}\frac1{1+\beta y_{u(i)}}
  \sum_{\gamma}\beta^{s(\gamma)}q^{D(\gamma)}\,
  \mathcal T_k\,\G^q_{\operatorname{end}(\gamma)},
\]
the sum over Left-free $(k-1)$-chains $\gamma$ from $u$ (one for each
Left-free endpoint, by Corollary~\ref{cor:rotation}). Expanding
$\mathcal T_k$ by \eqref{eq:line-operator}, the coefficient of
$\G^q_w$ on the right is a sum over \emph{composites} $(\gamma,\pi)$,
with $\pi$ a line path at $k$ from $\operatorname{end}(\gamma)$ to $w$,
of the weights
\[
 \beta^{s(\gamma)}q^{D(\gamma)}\,
 (-\beta)^{s(\pi)}\beta^{t(\pi)}q^{D(\pi)}\,
 \frac1{1+\beta y_{v(\pi)(k)}} .
\]
Since the $\G^q_w$ are linearly independent and $1-\beta^2q_{k-1}$ is
a nonzero polynomial, \eqref{eq:a-zero} at level $k$ is equivalent to
the identity, for every $w$,
\begin{equation}\label{eq:a-zero-junction}
\begin{split}
 &\sum_{(\gamma,\pi)\colon u\to w}
 \beta^{s(\gamma)}q^{D(\gamma)}
 (-\beta)^{s(\pi)}\beta^{t(\pi)}q^{D(\pi)}
 \frac1{1+\beta y_{v(\pi)(k)}}\\
 &\qquad=\frac{1-\beta^2q_{k-1}}{1+\beta y_{u(k)}}\,
  \beta^{s(u,w;k)}q^{D(u,w;k)}\,
  [\,w\in E(u,k)\text{ Left-free}\,],
\end{split}
\end{equation}
which is the junction identity \eqref{eq:junction-coeff} at $a=0$.
We prove \eqref{eq:a-zero-junction} by grouping composites by their
underlying path.

\emph{Words and cuts.}
A composite $(\gamma,\pi)$ determines the concatenated label sequence
$\omega=\gamma\pi$, which we call its \emph{word}; it is a directed
path from $u$ to $w$. By Lemma~\ref{lem:left-free} the labels of
$\gamma$ are $\prec$-increasing, so its labels with upper index $k$,
which have the largest possible $\prec$-value, form a terminal
segment of $\gamma$ with increasing lower indices; the left edges of
$\pi$ have upper index $k$ and decreasing lower indices. Hence
\[
 \omega=\gamma_{>k}\ (a_1',k)\,(a_2',k)\cdots(a_m',k)\ (k,b_1)\cdots(k,b_r),
\]
where $\gamma_{>k}$ consists of the labels of $\gamma$ with upper
index greater than $k$, the middle part is the \emph{block} of
$\omega$, and $b_1>\cdots>b_r>k$. The composite is recovered from
$\omega$ by the \emph{cut} $j\in\{0,\ldots,m\}$, the number of block
labels belonging to $\gamma$. Conversely, given a directed path
$\omega$ of this shape from $u$ to $w$, the cut $j$ yields a composite
exactly when $\gamma_{>k}(a_1',k)\cdots(a_j',k)$ is a Left-free
$(k-1)$-chain and $(a_{j+1}',k)\cdots(a_m',k)(k,b_1)\cdots(k,b_r)$ is
a line path, that is, by Lemma~\ref{lem:left-free} and
Definition~\ref{def:line-path}, exactly when
\[
 a_1'<a_2'<\cdots<a_j'
 \qquad\text{and}\qquad
 a_{j+1}'>a_{j+2}'>\cdots>a_m' .
\]
All composites with word $\omega$ have the same endpoint, the same
quantum degree $D(\omega)$, the same total number $|\omega|$ of
labels, and the same value $v(\pi)(k)$, namely the value $c(\omega)$
at position $k$ after the block, since position $k$ is touched only
by the block and the right edges. The weight of the composite with cut
$j$ is therefore
\[
 (-1)^{m-j}\,\beta^{|\omega|}q^{D(\omega)}\,\frac1{1+\beta y_{c(\omega)}} .
\]

\emph{Words with two valid cuts cancel.}
Let $m\geq1$. If $j$ is a valid cut with $1\leq j<m$, the largest
block entry is $a_j'$ or $a_{j+1}'$; if $j$ and $j'>j$ are both valid,
then $a_{j+1}',\ldots,a_{j'}'$ lie both in an increasing prefix and in
a decreasing suffix, so $j'=j+1$. Thus a word has at most two valid
cuts, and if it has two, they are consecutive and their weights
cancel. If $j$ is valid, then $j+1$ is valid exactly when
$a_j'<a_{j+1}'$, and $j-1$ is valid exactly when $a_j'>a_{j+1}'$
(for $1\leq j<m$; for $j=0$ or $j=m$ the block is monotone and the
cut $1$, respectively $m-1$, is valid as well). Hence a word with
exactly one valid cut $j$ has $1\leq j<m$ and $a_j'=a_{j+1}'$: the
same label $(a_j',k)$ is traversed twice in succession. Two
consecutive traversals of a label $(a,k)$ form a Bruhat edge followed
by a quantum edge or vice versa, which forces
$1=2(k-a)-1$, that is $a=k-1$. So a word with a single valid cut has
block
\begin{equation}\label{eq:single-cut-block}
 a_1'<\cdots<a_{j-1}'<k-1,\qquad a_j'=a_{j+1}'=k-1,\qquad
 k-1>a_{j+2}'>\cdots>a_m' .
\end{equation}

\emph{Moving the backtrack.}
From any permutation, $(k-1,k)$ is an edge (Bruhat if the value at
$k-1$ is smaller, quantum otherwise), and the pair
$(k-1,k)(k-1,k)$ returns to the starting permutation with quantum
degree $q_{k-1}$. Hence the pair can be inserted into or removed from
a word at any point without affecting the other edges. For a word
with block \eqref{eq:single-cut-block} and $m\geq3$, let $p$ be the
larger of the block entries adjacent to the pair ($a_{j-1}'$ or
$a_{j+2}'$, or the only one present), and move the pair to the other
side of $p$. The result is again a directed path to $w$ with a block
of the form \eqref{eq:single-cut-block}, with the same $|\omega|$,
$D(\omega)$ and $c(\omega)$, and its unique cut differs from $j$ by
one, so the two weights cancel. This is an involution without fixed
points on the words with a single valid cut and $m\geq3$.

\emph{Surviving words.}
What remains are the words with empty block,
$\omega=\gamma_{>k}(k,b_1)\cdots(k,b_r)$, and the words
$\omega=\gamma_{>k}(k-1,k)(k-1,k)(k,b_1)\cdots(k,b_r)$. Inserting or
removing the pair is a bijection between the two kinds, the second
kind having weight $-\beta^2q_{k-1}$ times the first (one more chain
label, one more left edge, degree $q_{k-1}$). In both kinds position
$k$ is untouched before the right edges, so $c(\omega)=u(k)$. The
left side of \eqref{eq:a-zero-junction} is therefore
\[
 \frac{1-\beta^2q_{k-1}}{1+\beta y_{u(k)}}
 \sum_{\omega}\beta^{|\omega|}q^{D(\omega)},
\]
the sum over directed paths $\omega=\gamma'(k,b_1)\cdots(k,b_r)$ from
$u$ to $w$ with $\gamma'$ a Left-free $(k-1)$-chain without upper
index $k$ and $b_1>\cdots>b_r>k$. By
Corollary~\ref{cor:sort-bijection} this sum equals
$\sum_\gamma\beta^{s(\gamma)}q^{D(\gamma)}$ over the Left-free
$k$-chains from $u$ to $w$, which by Corollary~\ref{cor:rotation} is
$\beta^{s(u,w;k)}q^{D(u,w;k)}$ if $w\in E(u,k)$ is Left-free and $0$
otherwise. This proves \eqref{eq:a-zero-junction}, hence
\eqref{eq:a-zero} at level $k$, and Theorem~\ref{thm:a-zero}. \qed

\begin{remark}
For $k=1$ the argument has no block and recovers
Proposition~\ref{prop:k-one} at $z=1/\beta$. The proof shows more
precisely that in the junction identity at $a=0$ the only composites
that do not cancel are the sorted composite
$\gamma_{>k}\,\big|\,(k,b_1)\cdots(k,b_r)$ of each Left-free endpoint
and its backtrack $\gamma_{>k}(k-1,k)\,\big|\,(k-1,k)(k,b_1)\cdots(k,b_r)$.
This is the pattern observed numerically for the moved case in
Section~\ref{sec:special} below; the general case differs in that
cut-sliding leaves residues proportional to $a$ (Lemma~\ref{lem:peel})
and in that chains with $\Left\neq\varnothing$ contribute.
\end{remark}

\begin{remark}[The specialization of Maeno--Naito--Sagaki]
A second one-point specialization of Theorem~\ref{prop:top} is in the
literature. Let
\[
 u=(n{-}k{+}1,\,n{-}k{+}2,\,\ldots,\,n{+}1,\ n{-}k,\,\ldots,\,2,\,1)\in S_{n+1},
\]
so that positions $1,\ldots,k+1$ carry the $k+1$ largest values in
increasing order and the remaining values follow in decreasing order.
The admissible $k$-chains from $u$ are exactly
\[
 \gamma_j=(k,k{+}1),\,(k{-}1,k{+}1),\,\ldots,\,(j,k{+}1),
 \qquad 1\leq j\leq k+1,
\]
with $\gamma_{k+1}$ the empty chain. Indeed, a first label $(a,b)$
with $b\geq k+2$ is neither Bruhat, as $u(a)>u(b)$, nor quantum, as
the value $n+1$ at position $k+1$ is not in $(u(b),u(a))$; so all
labels have upper index $k+1$, and after
$\gamma_j$ the only label $(a,k+1)$ with $a<j$ that is an edge is
$(j{-}1,k{+}1)$, because the value at $a$ is smaller than the value
$n-k+j$ at $k+1$, excluding a quantum edge, while the values at
positions $a+1,\ldots,j-1$ lie between them unless $a=j-1$. For $j\geq2$
position $1$ is untouched, so $u(1)\in\Fix$ and the coefficient
carries the factor $z+\overline y_{u(1)}=(1-a\,t_{u(1)})/(\beta t_{u(1)})$,
$t_c=1+\beta y_c$, which vanishes at $a=t_{u(1)}^{-1}$. For $j=1$ the
endpoint is
\[
 w=(n{-}k{+}2,\,\ldots,\,n{+}1,\ n{-}k{+}1,\ n{-}k,\,\ldots,\,1),
\]
with $\Left=\{u(2),\ldots,u(k)\}$, $\Out=\{u(1)\}$, $s=k$, $D=0$, and
coefficient $a^{k-1}t_{u(1)}^{-1}=t_{u(1)}^{-k}$ at $a=t_{u(1)}^{-1}$.
Since $\beta^kB_k^q=\sum_p(-a)^{k-p}F^k_p$, Theorem~\ref{prop:top}
predicts
\[
 \sum_{p=0}^k(-1)^{p}\,t_{u(1)}^{\,p}\,F^k_p\,\G^q_u=\G^q_w
 \qquad(\beta=-1),
\]
which is \cite[Proposition~2.3]{maeno2024presentation2} transported to
$QK_T(Fl_{n+1})$ as in their Section~3.1 (with $1-y_c=\mathbf e^{-\epsilon_c}$,
$u=w_{k+1}$ and $w=w_k$ in their notation), the key step in their
proof that quantum double Grothendieck polynomials represent Schubert
classes. Their proof of this special case expands each term of
$F^k_p$ by the general Chevalley formula for the weight
$w_0\epsilon_J$, $J\subseteq\{1,\ldots,k\}$.
\end{remark}

\subsection{The cancellation mechanism: cut composites and local moves}
\label{sec:mechanism}

Every term of \eqref{eq:junction-coeff} is a \emph{cut composite}:
a directed path in the quantum Bruhat graph
\[
 u=v_0\to v_1\to\cdots\to v_r=v\ \big|\ v\to\cdots\to w,
\]
whose first part is an admissible chain at level $k-1$ or $k-2$ and
whose second part is a line path at $k$; the bar marks the cut. The
term attached to it is the product of the chain coefficient
$C^{(j)}_{u,v}$ (times $a\beta q_{k-1}$ when $j=k-2$) and the line
weight $(-\beta)^{s}\beta^{t}q^{D}(1+\beta y_{v(\pi)(k)})^{-1}$.
The whole edge sequence from $u$ to $w$ determines $q^{D}$ for the
chain and the path together, and it determines $v(\pi)$; only the
position of the cut and the level distinguish the bookkeeping. The
cancellation in \eqref{eq:junction-coeff} is realized by two moves
on cut composites, which we now make precise, together with the
identity $\beta(z+\overline y_c)+a=(1+\beta y_c)^{-1}$ used in
\eqref{eq:diag-step}.

\begin{lemma}[Cut sliding]\label{lem:peel}
Let $\gamma$ be an admissible $(k-1)$-chain from $u$ to $v$, let
$\pi=(k-1,k)\cdot\pi'$ be a line path at $k$ from $v$ whose first
edge is $(k-1,k)$, and put $v'=vt_{k-1,k}$ and
$\kappa\in\{1,q_{k-1}\}$ for the weight of the edge $v\to v'$. Then:
\begin{enumerate}
\item $\gamma'=\gamma\cdot(k-1,k)$ is an admissible $(k-1)$-chain
 from $u$ to $v'$ if and only if $(k-1,k)\notin\gamma$, and $\pi'$
 is a line path at $k$ from $v'$ with $v(\pi')=v(\pi)$,
 $s(\pi')=s(\pi)-1$, $t(\pi')=t(\pi)$ and $q^{D(\pi)}=\kappa\,q^{D(\pi')}$.
\item If $v(k-1)\neq u(k-1)$ and $v(k)\neq u(k-1)$, then
 $\Fix,\Left,\Out$ computed for $v$ and for $v'$ agree, so that
 $C^{(k-1)}_{u,v'}=\beta\kappa\,C^{(k-1)}_{u,v}$, and the two terms of
 \eqref{eq:junction-coeff} indexed by $(v,\pi)$ and $(v',\pi')$
 cancel exactly.
\item If $v(k-1)=u(k-1)$, then $\gamma$ is also an admissible
 $(k-2)$-chain to $v$, $C^{(k-1)}_{u,v}=(z+\overline y_{u(k-1)})C^{(k-2)}_{u,v}$,
 and the two terms indexed by $(v,\pi)$ and $(v',\pi')$ sum to
 \begin{equation}\label{eq:peel-residue}
  a\kappa\,(-\beta)^{s(\pi')}\beta^{t(\pi')}q^{D(\pi')}\,
   \frac{C^{(k-2)}_{u,v}}{1+\beta y_{v(\pi)(k)}}.
 \end{equation}
\end{enumerate}
\end{lemma}

\begin{proof}
(1) The edge $(k-1,k)$ joins adjacent positions, so it is always an
edge of the quantum Bruhat graph, Bruhat if $v(k-1)<v(k)$ and quantum
with weight $q_{k-1}$ otherwise; it is the same edge in $\gamma'$ and
in $\pi$, whence the statement about weights. In $\gamma'$ the second
entries remain weakly decreasing since $k$ is the least possible
value, and the only condition of Definition~\ref{def:chain} that
involves the new last label is distinctness; condition~4 for the
previous label $(a_r,b_r)$ is automatic, because $b_r>k$, or $b_r=k$
and $a_r\neq k-1$. Removing the first edge of a line path leaves a
line path with the same left-part endpoint $v(\pi)$.

(2) Passing from $v$ to $v'$ exchanges the values $c=v(k-1)$ and
$d=v(k)$. The value $c$ leaves the window $\{1,\ldots,k-1\}$: it
cannot lie in $\Left$ (a value in position $k-1$ cannot have moved
left), and if it lies in $\Out$ or is not a window value of $u$ at
all, its status is unchanged. The value $d$ enters position $k-1$:
if $d=u(i)$ with $i\leq k-2$ it was in $\Out$ and remains there, and
otherwise it is not a window value. Hence $m$ is unchanged while the
chain length grows by one and the quantum weight is multiplied by
$\kappa$, giving $C^{(k-1)}_{u,v'}=\beta\kappa C^{(k-1)}_{u,v}$. The
line weights of $\pi$ and $\pi'$ differ by the factor $(-\beta)\kappa$,
so the two terms are negatives of each other.

(3) Here $v(k-1)=u(k-1)$, so by Corollary~\ref{cor:rotation}(2)
position $k-1$ is untouched by $\gamma$: no label of
$\gamma$ has first entry $k-1$, and $\gamma$ is admissible at level
$k-2$ with the same data; this gives the factorization of
$C^{(k-1)}_{u,v}$ as in \eqref{eq:diag-step}. In $v'$ the value
$c=u(k-1)$ has moved to position $k$, so it passes from $\Fix$ to
$\Out$ while $d$ behaves as in (2); thus
$C^{(k-1)}_{u,v'}=\kappa\,(1+\beta y_c)^{-1}(z+\overline y_c)^{-1}C^{(k-1)}_{u,v}$.
Writing $X=(-\beta)^{s(\pi')}\beta^{t(\pi')}q^{D(\pi')}(1+\beta y_{v(\pi)(k)})^{-1}$,
the term of $(v,\pi)$ is $X(-\beta)\kappa C^{(k-1)}_{u,v}$ and the term
of $(v',\pi')$ is $X\kappa(1+\beta y_c)^{-1}(z+\overline y_c)^{-1}C^{(k-1)}_{u,v}$.
Their sum is
\[
 X\kappa\,\frac{(1+\beta y_c)^{-1}-\beta(z+\overline y_c)}{z+\overline y_c}\,
   C^{(k-1)}_{u,v}
 =X\kappa a\,C^{(k-2)}_{u,v},
\]
which is \eqref{eq:peel-residue}.
\end{proof}

The case $v(k)=u(k-1)$, in which a fixed window value would enter
position $k-1$, does not occur in a cut composite whose endpoint
satisfies $w(k)=u(k)$: it did not appear in any of the verified
ranges below, and we do not use it. The residue
\eqref{eq:peel-residue} has exactly the shape of a level-$(k-2)$ term
of \eqref{eq:junction-coeff}, up to the factor $\beta q_{k-1}$: the
missing $q_{k-1}$ appears either through $\kappa=q_{k-1}$ or through a
factor $q_{k-1}$ inside $q^{D(\pi')}$, and the missing $\beta$ through
the line path $\pi'$ being one edge shorter. This is why the junction input
is $P_{k-1}(u)+a\beta q_{k-1}P_{k-2}(u)$ and not $P_{k-1}(u)$ alone.

\begin{lemma}[Braid move]\label{lem:braid}
Let $a<k-1$. For any permutation $v$, the two-edge sequences
$(k-1,k),(a,k)$ and $(a,k-1),(k-1,k)$ from $v$ have the same endpoint
$vt_{k-1,k}t_{ak}=vt_{a,k-1}t_{k-1,k}$, and if both are directed paths
in the quantum Bruhat graph they have the same quantum weight.
Consequently, if $\gamma$ is an admissible $(k-2)$-chain to $v$ with
$\gamma\cdot(a,k-1)$ admissible, and $\pi=(k-1,k)\cdot(a,k)\cdot\pi''$
is a line path from $v$, then $(k-1,k)\cdot\pi''$ is a line path from
$vt_{a,k-1}$ with the same $v(\pi)$, and the terms indexed by
$(v,\pi)$ and $(vt_{a,k-1},(k-1,k)\cdot\pi'')$ at level $k-2$ cancel
exactly whenever the window statistics of $v$ and $vt_{a,k-1}$
agree.
\end{lemma}

\begin{proof}
The identity $t_{k-1,k}t_{ak}=t_{a,k-1}t_{k-1,k}$ holds in $S_\infty$.
For the weights write $x=v(a)$, $y=v(k-1)$, $z=v(k)$. The adjacent
transposition $(k-1,k)$ is always an edge, quantum exactly when the
larger value is on the left. If $x<y$, the edge $(a,k)$ from
$vt_{k-1,k}$ is Bruhat, and its cover condition applied to the
intermediate position $k-1$, which holds $z$, gives $z\notin(x,y)$;
hence $[y>z]=[x>z]$, and both paths have weight $q_{k-1}^{[y>z]}$.
If $x>y$, the edge $(a,k)$ from $vt_{k-1,k}$ is quantum, so
$y<z<x$, and both paths have weight $q_aq_{a+1}\cdots q_{k-1}$.
The label $(a,k-1)$ is admissible at level $k-2$ since
$a\leq k-2<k-1$, and second entries remain weakly decreasing because
$k-1$ is the least possible value. The chain gains one edge
(a factor $\beta$ when $m$ is unchanged) while the line path loses
one left edge (a factor $-\beta$), and the quantum weights agree, so
the two terms cancel.
\end{proof}

When the window statistics change in Lemma~\ref{lem:braid} (the value
$v(a)$ or $v(k-1)$ is a fixed window value), a residue of the same
kind as \eqref{eq:peel-residue} is produced instead.

\paragraph{The drop term.}
The diagonal term of \eqref{eq:junction-coeff}, $\pi=\varnothing$ and
$v=w$, contributes $(1+\beta y_{w(k)})^{-1}\bigl(C^{(k-1)}_{u,w}+a\beta q_{k-1}C^{(k-2)}_{u,w}\bigr)$,
and after \eqref{eq:diag-step} what is left of it is the drop term
in \eqref{eq:diag-cancel}. In the language of cut composites it is
the composite with an empty line path, and Lemma~\ref{lem:peel}
pairs it with the composite
$u\xrightarrow{\gamma_w}w\xrightarrow{(k-1,k)}wt_{k-1,k}\ \big|\ (k-1,k)$,
for which the move of (1) is blocked because the label $(k-1,k)$ would
be repeated. Its term is $-\beta\kappa'\,(1+\beta y_{u(k)})^{-1}C^{(k-1)}_{u,wt_{k-1,k}}$ with
$\kappa\kappa'=q_{k-1}$, and by the computation in the proof of
Lemma~\ref{lem:peel} it equals
$-\beta^2q_{k-1}(1+\beta y_{u(k)})^{-1}C^{(k-1)}_{u,w}$ when
$w(k-1)\neq u(k-1)$, and
$-\beta q_{k-1}(1+\beta y_{u(k)})^{-1}(1+\beta y_{u(k-1)})^{-1}C^{(k-2)}_{u,w}
 =-\beta q_{k-1}(1+\beta y_{u(k)})^{-1}\bigl(\beta C^{(k-1)}_{u,w}+aC^{(k-2)}_{u,w}\bigr)$
when $w(k-1)=u(k-1)$, using \eqref{eq:diag-step} at level $k-1$. In
both cases it is the negative of the drop term, provided
$C^{(k-2)}_{u,w}=0$ in the first case.

\paragraph{A worked class.}
Take $u=21345$, $k=3$ and the endpoint $w=12345$, which is not in
$E(u,3)$, so the right side of \eqref{eq:junction-coeff} must vanish.
The nonzero terms are indexed by three cut composites of the single
edge sequence $(1,3),(2,3),(1,3)$ from $u$ to $w$ at level $2$, its
braid image at level $1$, and the drop:
\[
\begin{array}{lll}
 \text{level }2:& (1,3)\,\big|\,(2,3),(1,3) & \beta^2q_1q_2\,(z+\overline y_1)\,t_2^{-1}t_3^{-1},\\
 \text{level }2:& (1,3),(2,3)\,\big|\,(1,3) & -\beta q_1q_2\,t_1^{-1}t_2^{-1}t_3^{-1},\\
 \text{level }1:& (1,3)\,\big|\,(2,3),(1,3) & a\beta^3q_1q_2^2\,t_2^{-1}t_3^{-1},\\
 \text{level }1:& (1,3),(1,2)\,\big|\,(2,3) & -a\beta^3q_1q_2^2\,t_2^{-1}t_3^{-1},\\
 \text{drop}:& & a\beta q_1q_2\,t_2^{-1}t_3^{-1},
\end{array}
\]
with $t_c=1+\beta y_c$. The two level-$1$ terms are a braid pair
(Lemma~\ref{lem:braid}) and cancel. The two level-$2$ terms are a cut
sliding pair of type (3) in Lemma~\ref{lem:peel}, since the slid
value $v(2)=1=u(2)$ is fixed; their sum is the residue
$-a\beta q_1q_2t_2^{-1}t_3^{-1}$, which cancels the drop term. The
residue is absorbed by the drop rather than by a level-$1$ path term
because the relevant level-$1$ endpoint is $w$ itself.

\paragraph{Status.}
Lemmas~\ref{lem:peel} and~\ref{lem:braid} are proved. What remains
for the diagonal case is the global statement: the cut composites
with a given endpoint $w$, $w(k)=u(k)$, decompose into classes
generated by the two moves, and within each class the exact
cancellations and the residues of types \eqref{eq:peel-residue} sum
to zero against the level-$(k-2)$ terms and the drop. We have
verified this termwise, with the symbolic coefficients, for all
$u\in S_5$ and all $k$, and for a sample of $u\in S_6$ with
$k\leq5$; the identity \eqref{eq:diag-cancel} holds for every
endpoint $w$ in these ranges, including endpoints outside $E(u,k)$,
where it asserts the vanishing of the right side. A proof of the
global statement, and of its analogue for the moved cases below,
would complete the induction.

\paragraph{The moved cases.}
Suppose $u(k)\neq w(k)$. Then $C^{(k-1)}_{u,w}$ and
$C^{(k-2)}_{u,w}$ need not vanish, in contrast to the quantum
Schubert case: a $(k-1)$-chain may move position $k$ through a label
$(a,k)$. The identity \eqref{eq:junction-coeff} must then be
established by matching each $k$-chain to $w$ with a $(k-1)$- or
$(k-2)$-chain to some $v$ followed by a line path $v\to w$, and by
cancelling all remaining terms in pairs. The $k=1$ case shows the
pattern: there $P_{k-2}=0$, every nonempty line path is a single
up-path, and the matching is the identity. The complete case analysis
for $k\geq2$ is not carried out in this version.

\subsection{Special cases}
\label{sec:special}

Two specializations isolate the two difficulties of the junction
identity from each other.

\paragraph{The classical case \texorpdfstring{$q=0$}{q=0}.}
Set all $q_j=0$. Then all edges are Bruhat covers, the level-$(k-2)$
part of the junction input and the drop term vanish (both carry
$q_{k-1}$), and \eqref{eq:junction} reduces to
\begin{equation}\label{eq:junction-classical}
 \beta\,P_k(u)=\mathcal T_kP_{k-1}(u)-a\,P_{k-1}(u),
\end{equation}
the expansion of $\beta B_k+aB_{k-1}=(1+\beta x_k)B_{k-1}$ by the
classical equivariant $K$-Chevalley formula. The diagonal case is
then vacuous:

\begin{lemma}\label{lem:classical-diagonal}
Let $q=0$ and $w(k)=u(k)$. There is no admissible $(k-1)$-chain
$\gamma$ from $u$ to some $v$ together with a nonempty line path at
$k$ from $v$ to $w$. Consequently \eqref{eq:junction-classical} holds
at every such $w$, by \eqref{eq:diag-step}.
\end{lemma}

\begin{proof}
Follow the value in position $k$ along the composite path. Every
label $(a,k)$, whether in $\gamma$ or in the left part of the line
path, is a Bruhat cover, so it replaces the value in position $k$ by
a strictly smaller one; every right edge $(k,b)$ replaces it by a
strictly larger value coming from a position $b>k$. Suppose first
that some label $(a,k)$ occurs, and let $(\alpha,k)$ be the first.
It moves $u(k)$ to position $\alpha$, which is touched afterwards
only by a left edge $(\alpha,k)$ of the line path; that edge would
have to be a Bruhat cover with $u(k)=v(\alpha)<v(k)$, whereas every
value occupying position $k$ after $(\alpha,k)$ is smaller than
$u(k)$. Hence $u(k)$ never returns to position $k$, and the right
edges, which only import values from positions beyond $k$, cannot
return it either. If no label $(a,k)$ occurs, the line path consists
of right edges only, and the first of them moves $u(k)$ to a position
$b>k$ that is never touched again. In all cases $w(k)\neq u(k)$.
\end{proof}

For the moved case at $q=0$ the computations of
Section~\ref{sec:mechanism} show that only cut sliding is needed,
together with the end braid $t_{k-1,k}t_{kb}=t_{k-1,b}t_{k-1,k}$,
which converts a right edge following a junction edge $(k-1,k)$ into
two $(k-1)$-chain labels; the residues \eqref{eq:peel-residue} cancel
the term $-aC^{(k-1)}_{u,w}$ of \eqref{eq:junction-coeff} exactly as
in the example $u=21435$, $k=3$, $w=24315$, where the pair
$u\,\big|\,(2,3),(3,4)$ and $ut_{23}\,\big|\,(3,4)$ leaves the residue
$a\beta(z+\overline y_2)t_1^{-1}=aC^{(2)}_{u,w}$.

The word-and-cut bookkeeping of Section~\ref{sec:a-zero} organizes
the classical case completely at the level of individual words. We
normalize the coefficients by
\[
 \widehat C^{(j)}_{u,v}:=\beta^{j}C^{(j)}_{u,v}
 =\beta^{s}\prod_{i=1}^{j}\phi_i(u,v;j),
 \qquad
 \phi_i=\begin{cases}
  t_{u(i)}^{-1}-a& u(i)\in\Fix,\\
  a& u(i)\in\Left,\\
  t_{u(i)}^{-1}& u(i)\in\Out,
 \end{cases}
\]
with $t_c=1+\beta y_c$; this is \eqref{eq:coefficient} at $q=0$ written
with $z+\overline y_c=(t_c^{-1}-a)/\beta$ and $m+|\Fix|=j$. In this
notation the classical junction identity \eqref{eq:junction-classical}
at the endpoint $w$ reads
\begin{equation}\label{eq:classical-junction-hat}
 \widehat C^{(k)}_{u,w}+a\,\widehat C^{(k-1)}_{u,w}
 =\sum_{(\gamma,\pi)\colon u\to w}
   \beta^{s(\gamma)}(-\beta)^{s(\pi)}\beta^{t(\pi)}\,
   t_{v(\pi)(k)}^{-1}\prod_{i<k}\phi_i(u,\operatorname{end}\gamma;k-1),
\end{equation}
the sum over composites of an admissible $(k-1)$-chain $\gamma$ and a
line path $\pi$ of Bruhat edges. As in Section~\ref{sec:a-zero}, a
composite has a word
$\omega=\gamma_{>k}\,(a_1',k)\cdots(a_m',k)\,(k,b_1)\cdots(k,b_r)$ and a
cut $j$; now the chain part may contain descents, so the valid cuts
are no longer determined by monotonicity alone. Call the block entry
$a_r'$ \emph{fresh} if position $a_r'$ is touched neither by
$\gamma_{>k}$ nor by $a_1',\ldots,a_{r-1}'$, and write
$\Sigma(\omega)$ for the sum of the weights of all composites with
word $\omega$.

\begin{proposition}[Classical word lemma]\label{prop:classical-word}
Let $q=0$ and let $\omega$ be the word of a composite, with block
$(a_1',\ldots,a_m')$, and let $S$ be the length of the longest
strictly decreasing suffix $a_{m-S+1}'>\cdots>a_m'$ of the block
($S=0$ if $m=0$). Then:
\begin{enumerate}
\item The block entries are pairwise distinct.
\item The valid cuts form the interval $m-S\leq j\leq J$, where $J$ is
 the least index $r<m$ such that $a_r'$ is not fresh and
 $a_r'>a_{r+1}'$, or $J=m$ if there is no such $r$.
\item $\Sigma(\omega)\neq0$ if and only if $J=m$ and all entries
 $a_{m-S+1}',\ldots,a_m'$ of the decreasing suffix are fresh. In that
 case
 \[
  \Sigma(\omega)
  =\beta^{|\omega|}\,t_{c(\omega)}^{-1}\,a^{S}
   \prod_{i\notin\{a_1',\ldots,a_m'\}}\phi_i^{0}
   \prod_{r\leq m-S}\psi_r ,
 \]
 where $c(\omega)$ is the value at position $k$ after the block,
 $\phi_i^0=\phi_i(u,u\gamma_{>k};k-1)$, and for $r\leq m-S$:
 $\psi_r=\phi^0_{a_r'}$ if $a_r'$ is not fresh, $\psi_r=a$ if $a_r'$ is
 fresh and $a_r'>a_{r+1}'$, and $\psi_r=t_{u(a_r')}^{-1}$ if $a_r'$ is
 fresh and either $a_r'<a_{r+1}'$ or $r=m-S$.
\end{enumerate}
\end{proposition}

\begin{proof}
(1) A Bruhat edge $(a,k)$ requires the value at $a$ to be smaller
than the value at $k$, and it moves the smaller value to $k$; so along
the block the value at position $k$ strictly decreases. If $(a,k)$
occurred at two block indices $r<r'$, the value placed at $a$ by the
first occurrence would still be there at the second, and it exceeds
every later value at $k$; so the second occurrence is not a Bruhat
edge.

(2) By (1), an entry $a_r'$ is fresh exactly when position $a_r'$ is
not touched by $\gamma_{>k}$. The line part $(a_{j+1}',k)\cdots(a_m',k)(k,b_1)\cdots(k,b_r)$ is a
line path exactly when $a_{j+1}'>\cdots>a_m'$, that is $j\geq m-S$. The
chain part $\gamma_{>k}(a_1',k)\cdots(a_j',k)$ has distinct labels by
(1) and weakly decreasing upper indices, and condition~4 of
Definition~\ref{def:chain} is automatic inside $\gamma_{>k}$ and at
its last label. At a block index $r<j$ the lower index $a_r'$ is
repeated exactly when $a_r'$ is not fresh, since at level $k-1$ a
window position is touched only by labels with that lower index, and
then condition~4 requires $(a_r',k)\prec(a_{r+1}',k)$, that is
$a_r'<a_{r+1}'$. Hence the chain part is admissible exactly when
$j\leq J$.

(3) Fix a valid cut $j$ and let $v_j$ be the endpoint of the chain
part. We determine the status of each window position $i<k$ in $v_j$,
starting from its status in $v_0=u\gamma_{>k}$. Let $x_r$ be the value
at $a_r'$ just before $(a_r',k)$; after this label $x_r$ is at $k$, and
after $(a_{r+1}',k)$ (if $r<j$) it is at $a_{r+1}'$, where it stays,
since by (1) position $a_{r+1}'$ is not touched again by the block.
If $a_r'$ is fresh then $x_r=u(a_r')$, so position $a_r'$ has status
$\Out$ if $r=j$, status $\Left$ if $r<j$ and $a_{r+1}'<a_r'$, and status
$\Out$ if $r<j$ and $a_{r+1}'>a_r'$; before the label, for $r>j$, its
status is $\Fix$. If $a_r'$ is not fresh, then $x_r\neq u(a_r')$, and
we claim that the status of every value is unchanged by the block
labels of the chain part other than the fresh ones just described.
Indeed, a value $c=u(i)$ with status $\Left$ in $v_0$ sits at a position
$p<i$ reached by a descent $(i,b),(p,b)$ of $\gamma_{>k}$, and a chain
part containing $(p,k)$ would contain the subsequence
$(i,b),(p,b),(p,k)$ with $p<i\leq k-1<k<b$, which the prohibition
lemma \cite[Lemma~2.9]{lenart2007pieri} excludes for an admissible
chain. A value with status $\Out$ at a window position $p$ is moved
by $(p,k)$ to position $k$, and by the next label, if any, to
$a_{r+1}'$ with $a_{r+1}'>p$ by condition~4, since $p$ is a repeated
lower index; so it keeps status $\Out$. Values at positions
greater than $k$ are not moved by the block. This proves that the
weight of the composite with cut $j$ is
\[
 \beta^{|\omega|}(-1)^{m-j}\,t_{c(\omega)}^{-1}
 \prod_{i\notin\{a_r'\}}\phi_i^0
 \prod_{a_r'\text{ not fresh}}\phi^0_{a_r'}
 \prod_{a_r'\text{ fresh}}\phi_r^{(j)},
\]
where
\[
 \phi_r^{(j)}=\begin{cases}
  a& r<j,\ a_{r+1}'<a_r',\\
  t_{u(a_r')}^{-1}& r<j,\ a_{r+1}'>a_r',\ \text{or } r=j,\\
  t_{u(a_r')}^{-1}-a& r>j,
 \end{cases}
\]
the sign coming from the $m-j$ left edges of the line part.

Now sum over $j\in[m-S,J]$. If $J<m-S$ there is no valid cut and
$\Sigma(\omega)=0$; note that $J\neq m-S$ when $m-S\geq1$, since
$a_{m-S}'$ is an ascent by the maximality of $S$. For $r<m-S$ we have
$r<j$ throughout, so
$\phi^{(j)}_r=\psi_r$ is independent of $j$; for $r=m-S\geq1$, the
entry $a_{m-S}'$ is an ascent, so
$\phi^{(j)}_{m-S}=t_{u(a_{m-S}')}^{-1}$ both for $j=m-S$ and for
$j>m-S$. If $m-S<J<m$, the entries $a_r'$ with $m-S<r<J$ are fresh
descents, $a_J'$ is not fresh, and the factors for $r>J$ do not depend
on $j$. Writing $\tau_i=t_{u(a_{m-S+i}')}^{-1}$ and
$T_l=a^{l}\prod_{l<i<L+1}(\tau_i-a)$ for $0\leq l\leq L:=J-(m-S)-1$,
the remaining sum is, up to a common factor,
\begin{multline*}
 \sum_{j=m-S}^{J}(-1)^{m-j}\prod_{m-S<r<J}\phi_r^{(j)}\\
 =\pm\Bigl(T_0-(T_0+T_1)+(T_1+T_2)-\cdots
     +(-1)^{L}(T_{L-1}+T_L)+(-1)^{L+1}T_L\Bigr)=0,
\end{multline*}
because the cut $j=m-S+l$ with $1\leq l\leq L$ contributes
$a^{l-1}\tau_l\prod_{i>l}(\tau_i-a)=T_{l-1}+T_l$, the cut $j=m-S$
contributes $T_0$, and the cut $j=J$ contributes $a^L=T_L$. If $J=m$
and some suffix entry is not fresh, it must be $a_m'$ (the others are
descents), its factor $\phi^0_{a_m'}$ is constant, and the same
telescoping with $L=S-1$ gives $0$. If $J=m$ and all suffix entries
are fresh, the same computation with $L=S$, where now the cut $j=m$
contributes $a^{S-1}\tau_S=T_{S-1}+T_S$, gives $T_S=a^S$. Collecting
the constant factors yields the stated formula.
\end{proof}

Thus at $q=0$ a word contributes either nothing or a single product.
In the surviving words every entry of the decreasing suffix of the
block is a fresh position whose value, after the block, has moved to
the left (status $\Left$ at level $k$), and the suffix contributes
exactly the factor $a^S$ of \eqref{eq:coefficient} for these values.
What remains for the classical case is the matching: for every $w$,
the words with $\Sigma(\omega)\neq0$ and endpoint $w$ must have total
weight $\widehat C^{(k)}_{u,w}+a\widehat C^{(k-1)}_{u,w}$. We have
verified for all $u\in S_5$ and all $k\leq5$ that there is in fact
\emph{at most one} such word for each $w$, and that its $\Sigma(\omega)$
equals the right side; in the examples the surviving word is obtained
from the $k$-chain, or from the $(k-1)$-chain when $w\in E(u,k-1)$,
by commutations of disjoint transpositions as in
Lemma~\ref{lem:commute} and braid relations such as
$t_{kb}t_{ab}=t_{ak}t_{kb}$ and $t_{ab}t_{ak}=t_{ak}t_{kb}$, which
convert labels into junction labels $(a,k)$ and right edges $(k,b)$.
We have not proved this matching.

There is also a second route in the classical case, which has no
quantum analogue. Since $\prod_{j\leq k}(x_j+z)$ is symmetric in
$x_i,x_{i+1}$ for $i\neq k$, the isobaric divided difference
$\pi_i^x$ commutes with multiplication by it, and
$\pi^x_i\G_u=\G_{us_i}$ when $us_i<u$, $\pi^x_i\G_w=-\beta\G_w$ when
$ws_i>w$. Hence the classical coefficients must satisfy
\begin{equation}\label{eq:x-recursion}
 C_{us_i,w'}
 =\begin{cases}
   C_{u,w's_i}-\beta\,C_{u,w'}& w's_i>w',\\
   0& w's_i<w',
  \end{cases}
 \qquad(us_i<u,\ i\neq k,\ q=0).
\end{equation}
We have verified that the closed form \eqref{eq:coefficient} at $q=0$
satisfies \eqref{eq:x-recursion} for all $u\in S_4$ and all $k$ (1200
instances). The identity is suggestive: when a window value $c$ moves
from position $i+1$ to position $i$, the two terms on the right
combine as $(1+\beta y_c)^{-1}-\beta(z+\overline y_c)=1-\beta z$,
which is the origin of the factor $(1-\beta z)^{|\Left|}$. Granting
\eqref{eq:x-recursion} for the chain combinatorics, the classical
theorem for all $u$ follows by descending induction in the right weak
order from the case of $k$-Grassmannian $u$ (no descents outside
position $k$). For such $u$ the product is symmetric in
$x_1,\ldots,x_k$ and in the remaining variables, all endpoints are
again $k$-Grassmannian (checked for all $u\in S_6$), and the statement
is an equivariant $K$-theoretic Pieri rule for the Grassmannian,
multiplying by $\prod_{j\le k}(x_j+z)=(1-\beta z)^k\prod_j(x_j\oplus z')$
with $z'=z/(1-\beta z)$; the expansion of this product into the
classes $\G_{c[k,p]}$ is the case $u=e$ of Theorem~\ref{prop:top}.

\paragraph{The case \texorpdfstring{$k=2$}{k=2}.}
For $k=2$ a line path has at most one left edge, namely $(1,2)$, and
the level-$0$ part of the junction input is $a\beta q_1\G_u$. The
cut composites are the sequences
\[
 (1,b_1),\ldots,(1,b_r)\ \big|\ [(1,2)]\ (2,c_1),\ldots,(2,c_t),
\]
with the junction edge $(1,2)$ possibly absent, present on one side of
the cut, or present on both sides. The last possibility is a
\emph{backtrack}: a chain ending with a quantum edge $(1,2)$ followed
by the same edge reversed as a Bruhat line edge, or vice versa. Its
contribution relative to the composite without the two edges is
$\beta\cdot(-\beta)\cdot q_1=-\beta^2q_1$, and this is exactly how the
factor $1-\beta^2q_1$ on the left of \eqref{eq:junction} is
produced. For instance, for $u=21345$ and $w=32145\in E(u,2)$, the
only terms are $u t_{13}\,\big|\,(2,3)$, with value $\beta t_1^{-1}t_2^{-1}$,
and the backtrack $u t_{13}t_{12}\,\big|\,(1,2),(2,3)$, with value
$-\beta^3q_1t_1^{-1}t_2^{-1}$, whose sum is
$\beta(1-\beta^2q_1)C^{(2)}_{u,w}$. For $w=23145$ the backtrack
$ut_{12}\,\big|\,(1,2),(2,3)$ has value $-\beta^2q_1t_1^{-1}t_2^{-1}$ and the
level-$0$ composite $u\,\big|\,(2,3)$ has value $a\beta^2q_1t_1^{-1}$;
their sum is $-\beta^3q_1(z+\overline y_2)t_1^{-1}=-\beta^2q_1\,\beta C^{(2)}_{u,w}$,
again by $(1+\beta y_2)^{-1}-a=\beta(z+\overline y_2)$. Thus for
$k=2$ all three mechanisms---cut sliding, backtracks, and the
absorption of residues by the level-$0$ terms---are visible, while
the braid move of Lemma~\ref{lem:braid} cannot occur. We have
verified \eqref{eq:junction-coeff} for $k=2$ termwise for all
$u\in S_5$. The diagonal case can be settled completely.

\begin{proposition}[Diagonal case for $k=2$]\label{prop:k2-diag}
Let $k=2$ and $w(2)=u(2)$. Then \eqref{eq:diag-cancel} holds; hence
\eqref{eq:junction-coeff} holds at $w$.
\end{proposition}

\begin{proof}
We classify the nonempty cut composites ending at $w$. Position $2$
is touched by every edge of a line path at $2$. A right edge $(2,c)$
sends the value in position $2$ to the position $c>2$, which is never
touched again, so if the line path contains a right edge, the final
value in position $2$ is not $u(2)$ unless $u(2)$ was brought to
position $c$ earlier---impossible, as $c$ is touched only once. Hence
the line path is $\pi=(1,2)$, and the value $u(2)$ sits in position
$1$ of $v$, so the chain $\gamma$ to $v$ contains the label $(1,2)$.
For a level-$1$ chain this label is the last one, as $2$ is the least
admissible second entry, and $(1,2)$ moves $u(2)$ to position $1$ only
if it is the first label whose second entry is $2$; since labels are
distinct, $\gamma=\gamma_0\cdot(1,2)$ with $(1,2)\notin\gamma_0$, and
$w=vt_{12}=u\gamma_0$. Thus $\gamma_0$ is the unique $1$-chain from $u$
to $w$ (Proposition~\ref{prop:k-one}), and there is exactly one
level-$1$ composite, the backtrack $\gamma_w\cdot(1,2)\,\big|\,(1,2)$.
There is no level-$0$ composite, because a line path from $u$ cannot
restore $u(2)$ either. Write $\kappa,\kappa'$ for the weights of the
two traversals of $(1,2)$, so $\kappa\kappa'=q_1$.

If $w\neq u$, the window value $u(1)$ is in $\Out$ for both $w$ and
$wt_{12}$ (it has left position $1$ and $w(1)\neq u(1)$), so by
Lemma~\ref{lem:peel}(2) $C^{(1)}_{u,wt_{12}}=\beta\kappa C^{(1)}_{u,w}$,
and the backtrack contributes
$(-\beta)\kappa'(1+\beta y_{u(2)})^{-1}\beta\kappa C^{(1)}_{u,w}
=-\beta^2q_1(1+\beta y_{u(2)})^{-1}C^{(1)}_{u,w}$,
while the drop term is $\beta q_1(1+\beta y_{u(2)})^{-1}\beta C^{(1)}_{u,w}$
because $C^{(0)}_{u,w}=0$. They cancel.

If $w=u$, then $\gamma_0$ is empty, $C^{(1)}_{u,u}=z+\overline y_{u(1)}$,
and $u(1)$ passes from $\Fix$ to $\Out$ in $ut_{12}$, so
$C^{(1)}_{u,ut_{12}}=\kappa(1+\beta y_{u(1)})^{-1}$ and the backtrack
contributes $-\beta q_1(1+\beta y_{u(1)})^{-1}(1+\beta y_{u(2)})^{-1}$.
The drop term is
$\beta q_1(1+\beta y_{u(2)})^{-1}\bigl(\beta(z+\overline y_{u(1)})+a\bigr)
=\beta q_1(1+\beta y_{u(1)})^{-1}(1+\beta y_{u(2)})^{-1}$. They cancel.
\end{proof}

The moved case for $k=2$ requires in addition the classification of
$2$-chains: their labels form segments $(1,b)$, $(2,b)$, $(1,b)(2,b)$
or $(2,b)(1,b)$ for strictly decreasing $b$, the last form being
allowed only at the first occurrence of the lower index $2$. The
computations through $S_5$ exhibit the following pattern, which is
what a proof must establish. If $w\in E(u,2)$ and its $2$-chain has
no segment of the form $(2,b)(1,b)$, then sorting the chain, all
labels $(1,b)$ first and all labels $(2,c)$ afterwards, gives a cut
composite $\gamma^\ast\,\big|\,\pi^\ast$ whose term is
$\beta C^{(2)}_{u,w}$, and its backtrack
$\gamma^\ast(1,2)\,\big|\,(1,2)\pi^\ast$ contributes
$-\beta^3q_1C^{(2)}_{u,w}$; for instance $u=24135$, $w=31524$, with
chain $(2,5),(1,4),(2,3)$ and composite $(1,4)\,\big|\,(2,5),(2,3)$.
A segment $(2,b)(1,b)$ is where a window value moves left, and it is
resolved by $t_{2b}t_{1b}=t_{1b}t_{12}$: for $u=e$ and $w=231$, with
chain $(2,3),(1,3)$, the composites are $(1,2)\,\big|\,(2,3)$ with term
$\beta t_1^{-2}$ and $\varnothing\,\big|\,(1,2),(2,3)$ with term
$-\beta^2(z+\overline y_1)t_1^{-1}$, a cut sliding pair of type (3)
whose residue $\beta a t_1^{-1}$ is exactly $\beta C^{(2)}_{e,231}$:
the factor $1-\beta z$ attached to a value in $\Left$ is produced by
$(1+\beta y_c)^{-1}-\beta(z+\overline y_c)=1-\beta z$. The level-$0$
composite $e\,\big|\,(1,2),(2,3)$, with term $-a\beta^3q_1t_1^{-1}$,
then cancels the part $\beta^3q_1C^{(2)}_{e,231}$ of the left side.
When the $1$-chain to $w$ ends with $(1,2)$, the term
$-a(1-\beta^2q_1)C^{(1)}_{u,w}$ is cancelled by the residue of the
pair $u\,\big|\,(1,2)\pi'$, $ut_{12}\,\big|\,\pi'$ together with the
level-$0$ composite $u\,\big|\,(1,2)\pi'$, using
$t_{12}t_{2c_1}\cdots t_{2c_t}=t_{1c_1}\cdots t_{1c_t}t_{12}$ and
Postnikov's lemma for the equality of weights. The outstanding point
is the validity in the quantum Bruhat graph of the sorted and braided
sequences; Lemma~\ref{lem:sort} settles the sorting for chains
without descents, and Lemma~\ref{lem:commute} is the local step for
moving a right edge past a label with smaller upper index in general,
but the braided segments $(2,b)(1,b)$ are not covered, and we have
not completed this case analysis.

\paragraph{The \texorpdfstring{$y$}{y}-recursion.}
Finally, there is a reduction that is valid for all $q$ and does not
use the Chevalley formula at all. The quantization map commutes with
the isobaric divided differences $\pi^y_i$ in the second alphabet,
and $\pi^y_i\G^q_w=\G^q_{s_iw}$ if $s_iw<w$, $\pi^y_i\G^q_w=-\beta\G^q_w$
otherwise. Applying $\pi^y_i$ to \eqref{eq:main} for $s_iu>u$ and
using $\pi_i(fg)=(1+\beta y_{i+1})\partial_i(f)g+s_i(f)\pi_i(g)$ shows
that the theorem for $s_iu$ implies the theorem for $u$ provided the
coefficients satisfy
\begin{equation}\label{eq:y-recursion}
 C_{u,w}=\pi^y_i\bigl(C_{s_iu,w}\bigr)
 +\begin{cases}
   s_i\bigl(C_{s_iu,\,s_iw}\bigr)& s_iw>w,\\
   \beta\,s_i\bigl(C_{s_iu,w}\bigr)& s_iw<w.
  \end{cases}
\end{equation}
We have verified \eqref{eq:y-recursion} for the closed form, with
$q$ arbitrary, for all $u\in S_5$ and $k\leq3$. It reduces the
theorem for all $u\in S_n$ to the single permutation $u=w_0^{(n)}$,
for which the chains are indexed by subsets of $\{1,\ldots,k\}$ in the
classical case; the quantum case of this base is not simpler than the
general one, which is why we rely on the Chevalley induction instead.

\section{A small example and specializations}

\begin{example}\label{ex:s1}
Take $u=21$ and $k=1$. The endpoints are $21,312,132$, and the
identity permutation $e$. Their chain data $(s,D)$ are respectively
\[
 (0,0),\qquad(1,0),\qquad(2,(1,0,\ldots)),\qquad
 (1,(1,0,\ldots)).
\]
For instance, the path to $132$ is
\[
 213\xrightarrow{(1,3)}312\xrightarrow{(1,2)}132,
\]
whose second edge is quantum. Since
$B_1^q=x_1+z-\beta q_1(1+\beta x_1)$,
the formula reads
\begin{equation}\label{eq:example}
\begin{split}
 \bigl(x_1+z-\beta q_1(1+\beta x_1)\bigr)\G_{21}^q
 ={}&
 \left(z-\frac{y_2}{1+\beta y_2}\right)\G_{21}^q\\
 &+\frac{\G_{312}^q+\beta q_1\G_{132}^q+q_1}
              {1+\beta y_2}.
\end{split}
\end{equation}
Here $\G_e^q=1$, and all polynomials have alphabets $(x;y)$.
This particular identity can be verified by expanding the four
quantum double Grothendieck polynomials and clearing the denominator.
\end{example}

\paragraph{Classical specialization.}
When all $q_i=0$, only paths without quantum edges contribute.
The multiplier becomes $\prod_{i=1}^k(x_i+z)$, giving the
classical double Grothendieck top-block formula as a specialization
of Theorem~\ref{prop:top}. In particular, the
equivariant classical specialization is itself new, rather than
an appeal to a previously published top-block formula.

\paragraph{Non-equivariant specialization.}
When all $y_a=0$, the coefficient becomes
\[
 q^D\beta^{s-m}z^{|\Fix|}(1-\beta z)^{|\Left|}.
\]
In addition setting $z=0$ leaves only endpoints with $\Fix=\varnothing$,
and the surviving coefficient is $q^D\beta^{s-k}$.
At $\beta=-1$ it is $(-1)^{s-k}q^D$.
This specialization is Theorem~\ref{thm:noneq}, derived from the quantum
Grothendieck Pieri formula of Naito--Sagaki \cite{naito2024pieri}
through the uniqueness of chains and Lemma~\ref{lem:h-left}.

\paragraph{Cohomological specialization.}
At $\beta=0$, terms with $s>m$ vanish. Those with $s=m$ have coefficient
\[
 q^D\prod_{a\in\Fix}(z-y_a).
\]
The corresponding comparison is with the top quantum double
Schubert block in coefficient alphabet $-y$.

\paragraph{Usual K-theory variables.}
At $\beta=-1$ the multiplicative equivariant variables in these
conventions may be written
\[
 Y_a=\frac1{1-y_a}.
\]
Thus comparisons with formulas in the $Y$ variables use
$y_a=1-Y_a^{-1}$. This change of variables does not alter the
permutation indices or the quantum parameters $q_i$.

\bibliographystyle{acm}
\bibliography{quantum_k_pieri}

\end{document}
